%-----------------------------------------------------------------------
%                          hires2_correction.tex
% A&A LaTeX class (aa.cls)
%-----------------------------------------------------------------------
%\documentclass[referee]{aa}    % for a referee version
\documentclass{aa}
\usepackage{graphicx}
\usepackage{txfonts}
\usepackage{placeins}
\usepackage{natbib}
\usepackage[colorlinks,allcolors=blue,bookmarks=false]{hyperref}
\usepackage{xcolor}
\usepackage{amsmath}
\hyphenation{get-sources get-filaments get-images non-uniform back-ground}

\newcommand{\draftnote}[1]{\textcolor{magenta}{\textbf{[#1]}}}
\newcommand{\chg}[1]{\textcolor{magenta}{#1}}
\newcommand{\mH}{m_{\rm H}}
\newcommand{\muH}{\mu_{\rm H_2}}
\newcommand{{\um}}{\,$\mu$m}
\newcommand{\Teff}{T_{\rm F}}
\newcommand{\Dimc}{\mathcal{D}_{\rm C}}
\newcommand{\Dim}{\mathcal{D}}
\newcommand{\Iim}{\mathcal{I}}
%%\newcommand{\Sig}{\Sigma}

\begin{document}

\title{Correcting the temperature bias of dust surface density images\\of molecular clouds}

\author{Alexander Men'shchikov\inst{1}}
\institute{D\'epartement d'Astrophysique (DAp), IRFU, CEA Saclay,
Universit\'e Paris-Saclay, 91191 Gif-sur-Yvette, France}

\date{Received / Accepted}
\offprints{Alexander Men'shchikov}
\titlerunning{Correcting the temperature bias of surface density images}

\abstract{Surface densities of molecular clouds are derived by fitting one modified blackbody to each pixel of far-infrared images.
In dense structures, the fit is dominated by the warmer material along the line of sight and underestimates the surface densities
by more than a factor of two in the densest material. The observed spectral shapes cannot reveal this underestimation, because its
signature corresponds to only a few percent of intensities. This paper presents a correction method for surface densities $\Sigma$,
calibrated on radiative transfer models. The dust temperature $T$ depends only on the effective attenuating column $\Sigma_{\rm A}$
that shields the dust from the interstellar radiation field. A single relation $T(\Sigma_{\rm A})$ describes $2086$ embedded models
of Bonnor--Ebert spheres and $7$ cylindrical filaments with a scatter of $4.6$\%. The correction factor is an explicit function of
$\Sigma$ alone and only slightly depends on the shapes of structures. The method applies to any conventional image of surface
density, obtained with any angular resolution. The method was tested on additional $1815$ models of cores, on isolated filaments,
and on a simulated sky built from radiative transfer models. The correction raises the median core mass from $0.63$ to $0.98$ of
the true mass, the linear densities of filaments from $0.60$--$0.96$ to $0.94$--$1.01$, and the total $\Sigma$ of the simulated sky
to $0.96$. The correction includes dependence on a wide range of strengths $G_{0}$ of the interstellar radiation field, measured on
$9$ supplementary grids of models. Applied to the molecular cloud in Aquila, the correction raises the total $\Sigma$ by a factor
of $1.89$. For the adopted dust properties, the total uncertainties of corrected surface densities are ${\sim}30$\% in diffuse
material and ${\sim}50$\% in dense structures, and those of corrected temperatures ${\sim}2$\,K.}

\keywords{ISM: clouds -- ISM: structure -- Stars: formation -- Infrared: ISM -- Submillimeter: ISM -- Techniques: image processing}

\maketitle
\nolinenumbers 

% =====================================================================

\section{Introduction}
\label{sec:intro}

Far-infrared and submillimeter imaging surveys of the nearby and distant molecular clouds with the \emph{Herschel} Space
Observatory \citep{Pilbratt_etal2010} have provided most of the current knowledge about the structures in star-forming regions
\citep[e.g.,][]{Andre+2010, Motte+2010, Andre+2014}. One of the most important quantities is the surface density $\Sigma$ of dusty
gas, derived by fitting an optically thin modified blackbody to the intensities $I_{\lambda}$ of each pixel \citep{Hildebrand1983}.
The fitting gives one dust temperature $T$ and one surface density $\Sigma$ for each line of sight.

Dense cores and filaments are extracted from such images, and every quantity derived from them inherits their accuracy.
Conventional surface density images have a poor angular resolution of $36.3{\arcsec}$ of the longest waveband at $500${\um}. A
multi-resolution differential scheme \citep{Palmeirim+2013} can enhance the resolution to $18.2{\arcsec}$. This work employs the
improved \emph{hires} method \citep{Menshchikov2021method} that produces sharper surface density images at the resolutions of all
observed wavebands. The bias discussed below affects both the conventional and high-resolution images in the same way, because they
are based on the same method of fitting pixel intensities.

A weak point of the fitting procedure is the assumption of a constant temperature. Although it is accurate for diffuse material,
whose temperatures change little along the line of sight, the assumption fails for a cold dense structure embedded in a warmer
cloud. The derived temperature is weighted by the emission, hence by the warmer material, which leads to the underestimated surface 
densities. The size of the effect on surface density follows from the sensitivity of the Planck function to the temperature,
\begin{equation}
   \frac{\partial \ln \Sigma}{\partial \ln T} = -\,\frac{x_\lambda}{1-{\exp\left(-x_\lambda\right)}}, \quad
   x_\lambda \equiv \frac{h_{\rm P} c}{\lambda k_{\rm B} T},
\label{eq:sensitivity}
\end{equation}
where $\lambda$ is the wavelength, $c$ is the speed of light, $h_{\rm P}$ is the Planck constant, and $k_{\rm B}$ is the Boltzmann constant. At $T
= 10$\,K, this derivative equals $-9.0$, $-5.8$, $-4.2$, and $-3.1$ at $160$, $250$, $350$, and $500${\um}, respectively. If a cold
core's temperature is overestimated by only $1$\,K, the surface density derived at $500${\um} is too low by $24$\% and the one
derived at $160${\um} is underestimated by a factor of $2.3$.

The bias cannot be absorbed into the assumed dust opacity, because it is not a constant factor. It grows with the surface density
of the surrounding cloud, because denser clouds produce larger temperature contrasts along the line of sight. The bias also grows
with the concentration of a structure, because a more concentrated core is better shielded. Two structures with the same true
$\Sigma$ in different environments are therefore recovered with different derived values.

The bias has been studied with radiative transfer models of clouds, cores, and filaments \citep{Shetty+2009b, Malinen+2011,
Juvela+Ysard2012, Juvela+2013}. These studies found that single-temperature fits underestimate the surface density toward the
densest structures, by up to a factor of $3$ \citep{Juvela+2013}. Several ways to recover the temperature structure were proposed.
The observed profiles of a core can be inverted under the assumption of spherical symmetry \citep{Roy+2014}. The emission of each
line of sight can be described by several temperature components and fitted in many wavebands at once \citep{Marsh+2015}. Radiative
transfer models can be fitted to the observed images of a whole field \citep{Juvela+2013}, and empirical models can be used to
combine the wavebands of different resolutions \citep{Juvela+Montillaud2013}. These methods need either an assumed geometry of a
particular structure or a model of the whole field. The correction presented here needs neither. It is calibrated once on radiative
transfer models and is then applied to any surface density image.

This paper presents a correction of the bias in the derived images that needs only the surface density that a single modified
blackbody fit returns. 
%The multi-resolution reconstruction with the \emph{hires} method is an important special case. We show that
%the images of a reconstruction at higher resolutions carry the same fractional deficit as its base lowest-resolution image.
%The correction factor can therefore be evaluated on the image being corrected.
Section~\ref{sec:models} describes the radiative transfer models used to calibrate and test the correction.
Section~\ref{sec:law} derives from the models the relation between the dust temperature and the column that shields the dust.
Section~\ref{sec:corr} turns this relation into the correction. Section~\ref{sec:valid} tests the correction on independent
radiative transfer models and on a simulated sky. Section~\ref{sec:aquila} applies it to the Aquila cloud observed with
\emph{Herschel}. Appendices~\ref{sec:limits}--\ref{app:tcorr} describe the accuracies of surface densities, masses, and
temperatures derived from \emph{Herschel} images and discusses the uncertainties, including the dependence on the radiation field
and the errors of the backgrounds of sources, and other details.

The following notation is used. Images are denoted by calligraphic letters, for instance $\Iim$ for an intensity image and $\Dim$
for a surface density image. The \emph{hires} surface density image at a high angular resolution of $13.5{\arcsec}$ is denoted as
$\Dim$ without a subscript, and $\Dim_{36}$ is the image derived with a $36.3{\arcsec}$ resolution. The surface density of a single
pixel or line of sight is $\Sigma$ as the number of hydrogen molecules per unit area, in cm$^{-2}$.

% =====================================================================

\section{Radiative transfer models}
\label{sec:models}

The surface density correction method is based on accurate radiative transfer models of spherical cores, cylindrical filaments, and
a diffuse cloud, in which the true surface density is known. All models were computed with the Monte Carlo dust continuum
radiative-transfer code \emph{radmc-3d}\footnote{\url{http://www.ita.uni-heidelberg.de/~dullemond/software/radmc-3d}}
\citep{Dullemond+2012} to obtain their one-, two-, or three-dimensional dust temperature structure and images of their surface
density and emergent surface brightness in the \emph{Herschel} photometric bands ($70$, $100$, $160$, $250$, $350$, and
$500$\,$\mu$m).

The models were heated from outside by the isotropic interstellar radiation field \citep{Black1994}, with a strength of $G_0=0.98$
in \citet{Habing1968} units. The radiation field was modified by an additional hot blackbody component with a temperature of
$2.2{\times}10^4$\,K to scale the field to higher values $G_0 > 1$. Integrating the far-ultraviolet part of the field over $4\pi$
steradian and comparing to the \citet{Habing1968} normalization gave the adopted fiducial value of $G_0 =
4\pi\,1.22{\times}10^{-3}/1.6{\times}10^{-3} \approx 10$. The dust opacity was that tabulated by \citet{Ossenkopf+Henning1994} for
grains with thin ice mantles, corresponding to coagulation time of $10^{5}$\,yr and hydrogen atom number density of
$10^{6}$\,cm$^{-3}$. The opacities were scaled by a factor of $1.0746$ to match the analytic approximation $\kappa_\lambda =
10\,(300\,\mu{\rm m}/\lambda)^2\,{\rm cm^2\,g^{-1}}$ commonly used in observational studies. The adopted fiducial distance of
$260$\,pc corresponds to the value previously used for the Aquila cloud \citep{Konyves+2015}. This section describes the three
structural components; more details are given in Appendices~\ref{sec:grid} and \ref{app:rtmodels}.

\subsection{Cores}
\label{sec:modcores}

Each model of a core adopted a radial volume density distribution $\rho(r)$ of an isothermal Bonnor--Ebert sphere \citep{Ebert1955,
Bonnor1956} embedded in a uniform spherical cloud. The main grid contained $2089$ models. The surface density of the embedding
cloud ranged from $3.75{\times}10^{20}$ to $9.6{\times}10^{22}$\,cm$^{-2}$ in steps of a factor of $2$. The mass of the
Bonnor--Ebert sphere ranged from $1.6{\times}10^{-3}$ to $31\,M_{\odot}$, and the half-maximum width of its volume density profile
from $0.003$ to $0.25$\,pc, both in logarithmic steps. Only the models whose peak surface density exceeded that of the embedding
cloud by at least $5$\% were computed. Their central-to-edge density ratios ranged from $1.04$ to $2.1{\times}10^{3}$. Three
models, with embedding clouds of $1.9$ to $7.9{\times}10^{23}$\,cm$^{-2}$, were optically thick in the far infrared. Optically thin
emission is required, so those models were not used. The relation of Sect.~\ref{sec:relation} (Eq.~(\ref{eq:law})) was fitted to
the other $2086$ models and to the models of filaments. A subdivision grid contained $1815$ models whose parameters lay between
those of the main grid. They were not used in the fit and served to test the correction (Sect.~\ref{sec:valgrid}). The additional
grid of radiation fields repeated $339$ models of the main grid, computed with $9$ field strengths, from $G_0 = 0.98$ to $483$, in
steps of a factor of $2.17$ (Sect.~\ref{sec:g0corr}).

%The main grid contained 2089 models, spanning wide ranges of the surface
%densities of the embedding cloud, the core sizes, and their central-to-edge volume density ratios. Three models reached the highest
%surface densities and were optically thick in the far infrared. 

\subsection{Filaments}
\label{sec:modfil}

Each filament was modeled as an isolated cylinder, without any surrounding cloud, with the length $L = 7.2$\,pc and computed on an
axisymmetric grid with the radial volume density profile
\begin{equation}
    \rho(r) = \rho_{\rm C}\left(1 + \left(2^{2/p}-1\right)\left(\frac{2r}{h}\right)^{2}\right)^{-p/2},
\label{eq:fdens}
\end{equation}
truncated at the radius $R = 0.5$\,pc, where $h$ is the half-maximum width of the volume density profile and the factor
$(2^{2/p}-1)$ makes $h$ independent of the slope $p$ \citep{Menshchikov2021method}. The profile resembles a Gaussian at small radii
($r \lesssim h$) and asymptotically approaches a power law $\rho(r)\propto r^{-p}$ at $r \gg h$. The width $h$ of each model was
chosen so that the surface density profile $\Sigma(r)$ of the truncated cylinder had the half-maximum width $H$ given below, which
is about twice $h$ for these shallow profiles. Only the middle section of each cylinder was used (far from its ends) that
approximates properties of an infinite cylinder. There, the radiation entering through the end faces of the cylinder is negligible
and the temperature depends only on the distance $r$ from the axis. Six of the eight filaments had $p = 1.75$, $h = 20.3{\arcsec}$,
$H = 40{\arcsec}$, and surface densities on the axis of $3{\times}10^{21}$, $10^{22}$, $3{\times}10^{22}$, $5{\times}10^{22}$,
$10^{23}$ and $2{\times}10^{23}$\,cm$^{-2}$. The other two filaments had $3{\times}10^{22}$\,cm$^{-2}$ on the axis, with $p = 2$
and $h = 23.7{\arcsec}$ ($H = 40{\arcsec}$) or with $p = 1.75$ and $h = 42.5{\arcsec}$ ($H = 80{\arcsec}$). The cylinder with
$2{\times}10^{23}$\,cm$^{-2}$ had an optical depth of $0.84$ at $100${\um} along the line of sight normal to its crest and it was
not used, like the three optically thick cores.

\begin{figure*}
\centering
\includegraphics[width=0.85\linewidth]{fig_geometry.pdf}
\caption{Geometry of the dust attenuation in different structures. \emph{Upper}: The radiation that heats a parcel of dust, marked
by the black point, arrives from every direction having crossed columns of very different lengths. The rays are drawn thicker where
less material has been crossed. Most energy arrives along the normal in the midplane of a layer; no direction is favored at the
center of a sphere; a cylindrical filament has intermediate properties -- resembling a layer along its axis and a sphere across the
axis. \emph{Lower left}: The effective attenuating column $\Sigma_{\rm A}$ of Eq.~(\ref{eq:seff}) against the fractional depth $u$
along the line of sight normal to a uniform-density layer (solid curves) and along the diameter of a uniform-density sphere (dashed
curves), for total surface densities $\Sigma$ of $10^{21}$, $10^{22}$ and $10^{23}$\,cm$^{-2}$ along that line of sight. The dotted
curves are the column along the shortest path to the surface, which for these central lines of sight is $\min(u, 1-u)\,\Sigma$ for
both shapes (for off-center lines of sight through a sphere or a filament, the nearest surface lies
across the line of sight, and the column is smaller), and the thin vertical line marks the mid-plane (the center of
the sphere), about which the profiles are symmetric. \emph{Lower right}: The dust temperature
that follows from $\Sigma_{\rm A}$ through Eq.~(\ref{eq:law}). The dotted line marks the floor temperature $T_{\rm min}$.}
\label{fig:geometry}
\end{figure*}

\subsection{Diffuse cloud and simulated sky}
\label{sec:modcloud}

Most of the observed area of a molecular cloud is covered by extended diffuse material, referred to here as the diffuse cloud, or
simply the cloud. Its was modeled as a rectangular box of $6.8{\times}6.8$\,pc with a $2$\,pc depth along the line of sight, whose
projection was a map with scale-free fluctuations, rising toward the center of the field. The map was derived from the same
scale-free background used previously in \citet{Menshchikov_etal2012,Menshchikov2021method,Menshchikov2021bench}. The volume
density was constructed to fluctuate along every line of sight with the same power spectrum as across the sky
(Appendix~\ref{app:rtmodels}). Its surface densities ranged from $3{\times}10^{20}$ to $3{\times}10^{22}$\,cm$^{-2}$, with a median
of $5{\times}10^{21}$\,cm$^{-2}$.

The simulated sky combined this cloud with seven filaments and $96$ cores in a $1793{\times}1793$ field with $3{\arcsec}$ pixels.
Every structural component was computed with the radiative transfer code, so that its emission and its surface density were
consistent. For the simulated sky, the same filaments were computed again, embedded inside cylindrical envelopes that were
constructed to shield them the way the the surrounding cloud does. The emission of their envelopes was subtracted, so that only the
filaments were added to the cloud. The filaments lay in the middle depth of the cloud, which was assumed to extend far in front of
and behind them. The cores were the grid models inside embedding spherical clouds whose surface density matched its surroundings in
the sky.

The emission of the simulated sky was convolved to the beams of the wavebands at $160$--$500${\um} and its surface density
$\mathcal{D}$ was derived with \emph{hires}, as done for real observations. The true surface density $\mathcal{D}_{\rm T}$ was also
convolved to the same $13.5${\arcsec} resolutions.

% =====================================================================

\section{What sets the temperature of the dust}
\label{sec:law}

\subsection{The radiation arrives from every direction}
\label{sec:why}

Dust in a molecular cloud is heated mostly by the interstellar radiation field, the ultraviolet and optical light of the stars of
the Galaxy that arrives from outside the cloud. Locally, it is also heated by stars that formed within the cloud, from protostars
to young clusters. A grain absorbs the radiation, warms up, and re-emits the energy in the far infrared. The temperature of a grain
is therefore set by the amount of radiation that reaches it. The models and the temperature bias correction of this paper consider
the external heating only, for which this amount is reduced by the material between the grain and the outside of the cloud. Internal
heating is discussed in Appendix~\ref{sec:notdo}.

The radiation arrives from all directions, and it crosses different amounts of material in each direction. Let $I_{\lambda 0}$ be
the unattenuated radiation field at a wavelength $\lambda$. Denoting $\Sigma_{\rm A}(\hat n)$ the column of hydrogen molecules
(cm$^{-2}$) that the radiation crosses on its way in along the direction $\hat n$, the mean intensity $J_{\lambda}$ that reaches
the grain is
\begin{equation}
  J_{\lambda} = \frac{1}{4\pi} \int I_{{\lambda}0}\,\exp\left(-\sigma_{\lambda}\,\Sigma_{\rm A}(\hat{n})\right){\rm d}\Omega,
\label{eq:meanint}
\end{equation}
where $\sigma_{\lambda}$ is the extinction cross section per hydrogen molecule. The quantity $\Sigma_{\rm A}(\hat n)$ is referred to
as the attenuating column for the incoming radiation.

The upper panels of Fig.~\ref{fig:geometry} illustrate this for the three shapes used to describe clouds. The rays are thicker
where the radiation has crossed less material, so they show which directions deliver the energy. Deep inside a layer, the radiation
arrives almost entirely along the normal. A direction inclined by an angle $i$ crosses $(\cos i)^{-1}$ times more material, and
grazing directions deliver nothing. At the center of a sphere, every direction crosses the same column, and no direction is
favored. A filament combines the two cases. Across the filament the path is short, and along its axis it is long. A filament
therefore resembles a sphere in its cross-section and a layer along its axis.

The mean intensity of Eq.~(\ref{eq:meanint}) is an integral over all directions, each with its own column $\Sigma_{\rm A}(\hat n)$.
It is convenient to describe the shielding of a parcel by a single number instead, the column that, if it were the same in all
directions, would give the same mean intensity. With a gray cross section $\sigma$ (see below), this effective attenuating column
$\Sigma_{\rm A}$ is defined by
\begin{equation}
  \Sigma_{\rm A} = - \sigma^{-1}
    \ln \left\langle\,\exp\left({-\sigma\,\Sigma_{\rm A}\!\left(\hat n\right)}\right)\,\right\rangle_{4\pi},
\label{eq:seff}
\end{equation}
where the angle brackets denote averaging over all directions. The exponential describes the fraction of the external radiation
that reaches the parcel from the direction $\hat n$. For an isotropic external field, their average is the ratio
$J_{\lambda}/I_{\lambda 0}$ of Eq.~(\ref{eq:meanint}), so that $\sigma\Sigma_{\rm A} = -\ln \left(J_{\lambda}/I_{\lambda 0}\right)$
is the optical depth that would reduce the radiation reaching the parcel by the same factor. 

This average in Eq.~(\ref{eq:seff}) is dominated by the least shielded directions. The effective column is therefore not the
arithmetic mean of $\Sigma_{\rm A}(\hat n)$ over directions, which is infinite for a layer, because grazing directions cross an
unlimited amount of material. It is larger than the column along the shortest path, from the parcel to the nearest point of the
surface, because the oblique directions deliver less radiation than the shortest path does. The dust temperature depends on
$\Sigma_{\rm A}$. Figure~\ref{fig:angular} in Appendix~\ref{app:shape} shows how the average of Eq.~(\ref{eq:seff}) is formed at
the mid-plane of a layer, and how $\Sigma_{\rm A}$ grows with the total surface density of the line of sight.

Two assumptions were made. First, the extinction at the heating wavelengths was taken to be gray, with $A_{V} =
\Sigma/(0.94{\times}10^{21}\,{\rm cm}^{-2})$ \citep{Bohlin+1978} and the optical depth $\tau = A_{V}/1.086 = \sigma\,\Sigma$, where
$\sigma = 9.8{\times}10^{-22}$\,cm$^{2}$. In reality, the extinction varies across the ultraviolet and optical, so
Eq.~(\ref{eq:seff}) is an average over wavelength as well as over direction. Second, the radiation field outside the cloud was
assumed to be isotropic, as in the models of Sect.~\ref{sec:models}.

\begin{figure}
\centering
\includegraphics[width=0.95\columnwidth]{fig_shielding.pdf}
\caption{The dust temperature against the effective attenuating column $\Sigma_{\rm A}$ in the radiative transfer models. The gray
scale is the mass-weighted density of the $102\,464$ shells of the $2086$ spheres. The solid line and the band are their weighted
median and their 16th to 84th percentiles in bins of $0.2$\,dex. The circles are the $315$ radii of the seven cylinders. The dashed
line is the relation of Eq.~(\ref{eq:law}), fitted to the spheres and the cylinders with equal weights. The dotted line marks the
temperature $T_{\rm min}$ that fully shielded dust approaches.}
\label{fig:shielding}
\end{figure}

\subsection{The relation}
\label{sec:relation}

The relation between $\Sigma_{\rm A}$ and $T$ is not a free choice. It follows from radiative transfer, once the dust properties
and the radiation field are fixed. The density structure of every model is known. The attenuating column of a shell was therefore
computed along any direction by integrating the density along that ray. Equation~(\ref{eq:seff}) was then evaluated by summing over
directions, and each shell was weighted by its mass. For each cylinder, $\Sigma_{\rm A}$ was computed for an infinite cylinder with
the same profile. Near the centers of the densest cores, few photon packets of the Monte Carlo calculation arrived, and the
temperatures of the innermost cells scattered widely. A shell among the ten innermost ones was rejected if its temperature differed
by more than $10$\% from the median of its five neighbors, together with every shell inside it. The rejected shells were $1.8$\% of
all shells, and they contained a negligible fraction of the mass. The fit does not depend on this choice.

The $102\,464$ shells of the $2086$ spheres and the $315$ radii of the seven cylinders follow a single relation
(Fig.~\ref{fig:shielding}). It was fitted with the function
\begin{equation}
T(\Sigma_{\rm A}) = T_{\rm min} + (T_0 - T_{\rm min}) \left(1 + \Sigma_{\rm A}/{\Sigma_\star}\right)^{-\gamma},
\label{eq:law}
\end{equation}
where $T_{\rm min} = 7.0454$\,K, $T_0 = 22.8302$\,K, $\Sigma_\star = 2.7317{\times}10^{21}$\,cm$^{-2}$ and $\gamma = 1.6828$. The
spheres and the cylinders were given equal total weights in the fit. The mass-weighted scatter about the relation is $4.6$\%. The
two limits have simple meanings. $T_0$ is the temperature of dust that the radiation reaches without attenuation. $T_{\rm min}$ is
the temperature that dust approaches when the radiation cannot reach it at all; the dust is then warmed only by the little
radiation that penetrates and by the microwave background. The spheres and the cylinders follow the same relation: on average, the
spheres are $0.8$\% warmer than the relation and the cylinders $0.3$ to $0.9$\% colder. Leaving out any one cylinder changes the
correction factor by at most $2$\%. Above $\Sigma_{\rm A} = 10^{23}$\,cm$^{-2}$, the models contain too little mass to constrain
the relation.

A single relation describes models that differ greatly from one another. This is the result that makes the method possible. The
size of a core and its concentration do not appear in Eq.~(\ref{eq:law}), and they do not need to. They affect the temperature only
by changing the amount of material between the dust and the radiation. The scatter about the relation, $4.6$\%, measures how much
structures of different shapes differ at the same $\Sigma_{\rm A}$. It is physical, and it would not decrease with more models.

% =====================================================================

\section{The correction}
\label{sec:corr}

\begin{figure*}
\centering
\includegraphics[width=0.76\textwidth]{fig_method.pdf}
\caption{The method step by step. The first group is the calibration, which is done once and yields the temperature relation of
Eq.~(\ref{eq:law}). The second is the calculation of Sect.~\ref{sec:factor}, repeated for every trial surface density, which
yields the table of the correction factor against the surface density that an observed image actually holds. Its last two steps
calibrate the low surface densities on a radiative transfer model of a diffuse cloud (Sect.~\ref{sec:diffuse}) and adapt the factor
to the radiation field of the cloud (Sect.~\ref{sec:g0corr}). The third group is
what is then done to an image, at any resolution.}
\label{fig:method}
\end{figure*}

\subsection{The calculation}
\label{sec:factor}

Apart from the calibration of diffuse material (Sect.~\ref{sec:diffuse}), nothing was left to adjust. For a line of sight with the
true surface density $\Sigma$, the factor was computed in five steps. (\emph{i}) For a structure of uniform density with the total
surface density $\Sigma$, the column $\Sigma_{\rm A}$ was computed at each depth from Eq.~(\ref{eq:seff}), once for a uniform sphere
along its diameter and once for a uniform plane-parallel layer along its normal. The two values were averaged at each depth, for
the reason given in Sect.~\ref{sec:orient}. (\emph{ii}) The temperature at each depth followed from Eq.~(\ref{eq:law}).
(\emph{iii}) The emergent intensity was obtained as the sum of the emission of all layers,
\begin{equation}
   \Iim_\lambda(\Sigma) = \kappa_\lambda\,\eta\,\muH\,\mH\,\Sigma
   \int_0^1 B_\lambda \big(T(\Sigma_{\rm A}(u))\big)\,{\rm d}u,
\label{eq:emission}
\end{equation}
where $u$ is the fraction of the total surface density that lies in front of a layer, $\kappa_\lambda$ is the dust opacity per unit
mass of dust, $\eta = 0.01$ is the dust-to-gas mass ratio, and $\mH$ is the mass of a hydrogen atom. (\emph{iv}) A single modified
blackbody was fitted to these intensities, in the same wavebands and with the same weights as the reconstruction. The fit gave a
temperature $\Teff$ and the surface density that the reconstruction would report, obtained from the reference waveband, which is
the longest one, $500${\um}:
\begin{equation}
   \Sigma_{\rm F}(\Sigma) = \frac{\Iim_{\rm ref}(\Sigma)} {\kappa_{\rm ref}\,\eta\,\muH\,\mH\,B_{\rm ref}(\Teff)},
\label{eq:fitted}
\end{equation}
(\emph{v}) The correction factor is the ratio of the two surface densities. It was tabulated against $\Sigma_{\rm F}$, because an
observed image contains $\Sigma_{\rm F}$ and not $\Sigma$:
\begin{equation}
   f(\Sigma_{\rm F}) = {\Sigma}/{\Sigma_{\rm F}} .
\label{eq:factor}
\end{equation}
Figure~\ref{fig:method} shows the five steps together with the steps before and after them.

\begin{figure}
\centering
\includegraphics[width=0.98\columnwidth]{fig_relation.pdf}
\caption{The correction factor against the surface density of the uncorrected image, for a uniform plane-parallel layer, for a
cube, for a uniform sphere, and for the average of the sphere and the layer. The dotted curve uses the column along the shortest
path instead of the angular average (Appendix~\ref{app:shape}). The thick magenta curve is the adopted relation for $G_0 = 10$.
Below $10^{21.6}$\,cm$^{-2}$ it replaces the analytic layer by a radiative transfer model of a diffuse cloud
(Sect.~\ref{sec:diffuse}), and above its maximum it is held constant (Sect.~\ref{sec:factor}). The thin magenta curves are the
adopted relation for $G_0 = 5$ and $33$, with the field factor of Eq.~(\ref{eq:factorg0}).}
\label{fig:relation}
\end{figure}

Figure~\ref{fig:relation} shows the result. For the analytic structures, the factor is unity below $10^{21}$\,cm$^{-2}$, where
nothing is attenuated. It rises steeply through the range of surface densities of most of a molecular cloud. It reaches its
maximum, $1.92$, near $3{\times}10^{22}$\,cm$^{-2}$ and then declines, to $1.47$ at $3{\times}10^{23}$\,cm$^{-2}$. At such surface
densities, most of the material of a uniform body has reached the temperature floor, so its line of sight is again nearly
isothermal. Real dense structures are not uniform. A concentrated core has a cold center inside a warmer envelope, and its
temperature still varies along the line of sight. The radiative transfer models confirm this: the factor that they require stays at
$2.0$ to $2.4$ up to $10^{23}$\,cm$^{-2}$, where the models end. The decline of the analytic relation is an artifact of the uniform
bodies. The factor was therefore held at its maximum above the surface density at which it is reached, so that it never decreases
with the surface density. The adopted relation differs from the analytic one in this respect and below $10^{22}$\,cm$^{-2}$, for
the reason given in Sect.~\ref{sec:diffuse}.

The lower left panel of Fig.~\ref{fig:geometry} shows the quantity that drives the calculation, the profile of $\Sigma_{\rm A}$
along a line of sight. At $10^{21}$\,cm$^{-2}$ the profile is nearly flat. The line of sight is then nearly isothermal, and there
is almost nothing to correct. As the surface density rises, the profile steepens and approaches the dotted line of the shortest
path. Once the material is optically thick, the oblique directions deliver nothing, and only the normal direction matters.


\subsection{Why the shape hardly matters}
\label{sec:orient}

The calculation needs a shape, but an observed field does not provide one. Cores are embedded in filaments. Filaments are curved
and blended with their neighbors. Everything lies on a background that fluctuates on all scales. It is therefore important that the
result should not depend much on the assumed shape. It does not.

Figure~\ref{fig:relation} gives the factor for a uniform sphere, a uniform cube and a uniform plane-parallel layer of the same
observed surface density. The three agree to within $10$\% everywhere. Outside the range from $10^{21.4}$ to
$10^{22.1}$\,cm$^{-2}$, they agree to within $4$\%. Figure~\ref{fig:geometry} shows why. At low surface densities nothing is
attenuated, and every shape gives $f = 1$. At high surface densities the deepest material is optically thick in every direction,
and the distribution of path lengths no longer matters. The shape matters only in between, where the correction itself is modest.
In short, where the geometry would matter, the correction is small, and where the correction is large, the geometry does not
matter. The average of the sphere and the layer was adopted, and half of their difference, at most $5$\%, was taken as the
uncertainty due to the geometry. The factor was computed for central lines of sight, along the normal of a layer and the diameter
of a sphere. Lines of sight through the outer parts of a compact structure pass closer to its surface. They are warmer and more
nearly isothermal than a central line of sight with the same $\Sigma$, and they are slightly over-corrected. The validation on all
pixels of the models includes this effect (Sect.~\ref{sec:valgrid}). Inclined layers and filaments are over-corrected
too (Appendix~\ref{sec:incl}).

Appendix~\ref{app:shape} explains why the sphere and the layer are the two limits. It also explains why the angular average cannot
be replaced by the column along the shortest path.

One assumption remains: each structure was assumed to have a uniform density. This assumption matters much less than it seems. For
a plane-parallel line of sight it does not matter at all. A parcel of dust is shielded by the column between the parcel and the
surface. This column does not depend on how the material within it is distributed. Putting half of the material into a dense clump
does not change the amount of material. The mass of a line of sight per unit column is also constant for any density profile,
because the column is the integral of the density. Let $\Sigma(s)$ be the column in front of a parcel at the depth $s$. The
emission of Eq.~(\ref{eq:emission}) integrated along the line of sight is then
\begin{equation}
   \int_0^L B_\lambda\big(T(\Sigma_{\rm A}(\Sigma(s)))\big)\,\rho(s)\,{\rm d}s
   = \int_0^{\Sigma} B_\lambda\big(T(\Sigma_{\rm A}(\Sigma^{\prime}))\big)\,{\rm d}\Sigma^{\prime} ,
\label{eq:invariance}
\end{equation}
and the density has disappeared. Two lines of sight with the same total surface density emit the same spectrum. It does not matter
whether their material is smooth or clumped, or whether a clump lies at the near side, in the middle, or at the far side.
Appendix~\ref{app:arrange} confirms this numerically for density profiles that differ by a factor of thirty. They reproduce the
smooth case to two parts in $10^{4}$.

Concentration and arrangement can therefore act only through the one thing that Eq.~(\ref{eq:invariance}) assumes: that the
radiation escapes along the line of sight and not across it. Their effect is modest. A sphere with a central-to-edge density ratio
of $10$ needs a correction $2$ to $4$\% larger than a uniform sphere of the same central surface density. For a ratio of $10^3$, it
needs $8$ to $19$\% more, and the difference grows with the surface density. A diffuse cloud with a denser structure embedded in it
is a different case. The correction for such a line of sight is between $13$\% too small and $23$\% too large, depending on the
depth of the structure in the diffuse cloud. Both effects are measured in Appendix~\ref{app:arrange}, and they enter the
uncertainties of Appendix~\ref{sec:limits}.

\subsection{Diffuse material}
\label{sec:diffuse}

The relation of Eq.~(\ref{eq:law}) was measured on spheres and cylinders. These are bounded structures. Most of the area of a
molecular cloud, however, is covered by the diffuse cloud (Sect.~\ref{sec:modcloud}). In each column of its radiative transfer
model, one modified blackbody was fitted to the emission at $160$ to $500${\um}, and the result was compared with the true surface
density.

The analytic layer, which enters the average of Sect.~\ref{sec:orient}, does not describe this cloud. The layer needs almost no
correction below $10^{21}$\,cm$^{-2}$, but the radiative transfer model of the cloud needs $5$\% even at
$5\times10^{20}$\,cm$^{-2}$, and $15$--$25$\% at a few $10^{21}$\,cm$^{-2}$. The reason is the surface of the cloud. Its outer
layers are heated by the unattenuated ultraviolet radiation from half the sky. They are up to $3$\,K warmer than Eq.~(\ref{eq:law})
predicts, because that relation describes the attenuation with a single gray opacity. The temperature therefore varies more along a
line of sight through the cloud than through the analytic layer, and the fit loses more.

The analytic layer was therefore replaced by the radiative transfer model of the cloud at low surface densities. Below a fitted
surface density of $10^{21.6}$\,cm$^{-2}$, the factor was set to the average of the analytic sphere and of the factor measured on
the cloud. Between $10^{21.6}$ and $10^{22}$\,cm$^{-2}$, it was blended smoothly into the analytic average of sphere and layer,
which was used above. The factor is thus $1.03$ at $3\times10^{20}$\,cm$^{-2}$ and $1.12$ at $3\times10^{21}$\,cm$^{-2}$, against
$1.00$ and $1.07$ for the analytic average. It remains a function of the surface density alone. The bounded spheres and the
extended cloud are the two limits for diffuse material, as the sphere and the layer are for a single structure. Their average
divides the uncertainty between them. A uniform cloud around a core is corrected about $2$\% too much (Sect.~\ref{sec:valgrid}).
The radiative transfer model of the cloud is corrected too little, by $6$\% at $10^{21}$\,cm$^{-2}$ and by $14$\% at
$10^{22}$\,cm$^{-2}$, with a median of $11$\% (Sect.~\ref{sec:sim}).

\subsection{The radiation field of a cloud}
\label{sec:g0corr}

The relation of Eq.~(\ref{eq:law}) and the factor $f$ were calibrated for $G_0 = 10$. Clouds are heated by very different fields.
The diffuse material of the regions listed in Appendix~\ref{app:fieldest} requires fields from below $1$ to about $250$. A stronger
field heats the dust at every attenuating column, and it changes the temperature contrast along the line of sight. The factor that
a line of sight requires therefore depends on the field (Appendix~\ref{sec:field}).

This dependence was measured on the grid of radiation fields (Sect.~\ref{sec:modcores}), in nine fields from $G_0 = 0.98$ to $483$.
For each field, the logarithm of the factor that each pixel of each model required was averaged in bins of $0.1$\,dex of the fitted
surface density $\Sigma_{\rm F}$. The ratio of this average to the same average at $G_0 = 10$ is the field factor $g(\Sigma_{\rm
F}, G_0)$. Above $10^{22.6}$\,cm$^{-2}$, the grid contains too few pixels to determine it, and $g$ was held at its value there.
Between the fields of the grid, $\log g$ is interpolated linearly in $\log G_0$. The field factor is within $3$\% of unity below
$10^{21.4}$\,cm$^{-2}$ in every field. At $10^{22}$\,cm$^{-2}$, it is $0.80$ for $G_0 = 1$, $0.91$ for $G_0 = 5$, $1.16$ for $G_0 =
33$, and $1.09$ for $G_0 = 250$ (Appendix~\ref{app:field}).

The factor applied to a cloud in the field $G_0$ is then
\begin{equation}
  f_{G_0}(\Sigma_{\rm F}) = \max_{\Sigma^{\prime} \le \Sigma_{\rm F}} \Big[ f(\Sigma^{\prime})\,g(\Sigma^{\prime}, G_0) \Big] ,
\label{eq:factorg0}
\end{equation}
which also holds the factor at its maximum (Sect.~\ref{sec:factor}). For $G_0 = 10$, $g = 1$. The field is a single number for a
cloud. It is the field that the models require to reproduce the observed temperatures, an effective field of the model dust, which
can differ from the physical field. It is estimated from the temperatures of the diffuse material of the whole map, with the
procedure of Appendix~\ref{app:fieldest}. It is not read from the temperature of each pixel, which is ambiguous
(Appendix~\ref{app:tcorr}). The correction therefore still depends on the surface density alone, and the field is an input of the
calculation, like the wavebands.

\subsection{How it is applied}
\label{sec:apply}

An image is corrected by multiplying it by the image of the correction factors evaluated on the image itself,
\begin{equation}
  \Dimc(O) = \mathcal{F}(O)\,\Dim(O), \quad \mathcal{F}(O) \equiv f_{G_0}\big(\Dim(O)\big) ,
\label{eq:correction}
\end{equation}
where $f_{G_0}$ of Eq.~(\ref{eq:factorg0}) is evaluated pixel by pixel. Without an estimate of the field, $G_0 = 10$ is used. This
applies to a conventional surface density map and to the \emph{hires} image at any resolution $O$. The factor was calibrated for the
surface density that a single modified blackbody fit returns at the resolution of the longest waveband. The \emph{hires} image at
$13.5{\arcsec}$ is sharper than any such fit, but it is also a surface density image, and the factor applies to it. Evaluated on
that image, the factor gives core masses closer to the truth than when it is evaluated on the image at $36.3{\arcsec}$, and on the
diffuse cloud the two agree to $0.5$\% (Sects.~\ref{sec:valgrid} and \ref{sec:sim}). Alternatively, the correction
$[f(\Dim_{36}) - 1]\,\Dim_{36}$ of the image at $36.3{\arcsec}$ can be added to the finer image, and Sect.~\ref{sec:base} compares
the two ways. The first was adopted, because it gives the sharper images and the more accurate masses. Correcting the intermediate
images of \emph{hires} before the reconstruction would require a separate calibration for each of them, because they are fitted to
fewer wavebands. The correction must be applied only once. The two implementations of the correction, in Fortran and in Python,
mark their output and refuse an input that has already been corrected. Both are distributed as part of \emph{getsf} on its home
page\footnote{\url{http://irfu.cea.fr/Pisp/alexander.menshchikov/}}.

\subsection{What the correction needs, and what it does not}
\label{sec:needs}

Nothing in Eqs.~(\ref{eq:seff}) to (\ref{eq:factor}) refers to the multi-resolution reconstruction. For a given field, the factor
is a function of one number: the surface density that a single modified blackbody fit returns at the resolution of the longest
waveband. This number can be obtained from the observed images alone. The base image of \emph{hires} is exactly this fit. This was
verified on the grid: a modified blackbody fitted to the same wavebands with the same weights reproduced the base image of
\emph{hires} to better than $1$\% for every model. The correction therefore applies to a conventional surface density map, made in
the usual way from the \emph{Herschel} intensities convolved to a common beam. Equation~(\ref{eq:correction}) then reduces to
multiplying that map by $f$.

What the \emph{hires} reconstruction adds, and what no correction can replace, is the information on scales finer than the beam of
the longest waveband, which it derives from the shorter wavebands. Without it, there is nothing to correct at $13.5{\arcsec}$, and
the corrected map is a map at $36.3{\arcsec}$.

The order of the operations matters when an image is decomposed into components. A parcel is shielded by the whole column along its
line of sight. The factor must therefore be evaluated on the total surface density, not on a component of it. The image must be
corrected first, and the background of a structure subtracted afterwards. If the background is subtracted first, the factor is
evaluated at a surface density that the material does not have, and the structure is corrected much too little. For the same
reason, the correction must never be applied to an image that is already corrected, or to a true surface density.

One condition applies when the correction is used outside \emph{hires}. The table of $f$ is built for a particular set of wavebands
and a particular weighting of them, because $\Sigma_{\rm F}$ depends on both. The table must be rebuilt if a map was made from
other wavebands or with other weights. The table cannot be replaced by anything derived from the photometry itself. The mixing of temperatures
along a line of sight broadens its spectral shape only slightly. At the centers of the cores (Sect.~\ref{sec:modcores}), a single
modified blackbody fits the four intensities to within $1$--$6$\%, while the fit loses $9$--$38$\% of $\Sigma$. A method that
derived the underestimation from these small residuals would amplify the photometric errors by factors of $7$ to $9$, far beyond the
calibration accuracy of \emph{Herschel} (Appendix~\ref{app:colors}). The correction must therefore come from models in which the
answer is known.

Two more properties follow from the construction. First, the correction is insensitive to noise, because it evaluates a smooth
function of an image and fits nothing. Gaussian noise of $1.5$, $0.50$, $0.30$, and $0.15$\,MJy\,sr$^{-1}$ at $160$, $250$, $350$,
and $500${\um} was added to the models. The total corrected surface density changed by less than $0.5$\%. Second, the correction
does not amplify the calibration errors of the wavebands much. The amplification by $7$ to $9$ found in Appendix~\ref{app:colors}
does not apply here, because the underestimation is not derived from the colors. A calibration error changes the fitted surface
density $\Sigma_{\rm F}$, and the corrected surface density $f(\Sigma_{\rm F})\,\Sigma_{\rm F}$ follows. The effect of an error of
$10$\% in each waveband was computed for uniform lines of sight of $3{\times}10^{21}$ to $10^{23}$\,cm$^{-2}$. An error at
$500${\um}, the reference waveband, changes the uncorrected surface density by $16$ to $18$\% and the corrected one by $16$ to
$23$\%. An error at $160$, $250$ or $350${\um} changes the fitted temperature by up to $0.4$\,K. It changes the uncorrected surface
density by $1$ to $12$\% and the corrected one by $1$ to $11$\%. An error of $10$\% in all wavebands at once changes the
uncorrected surface density by $10$\% and the corrected one by $9$ to $13$\%. Where $f$ rises steeply with $\Sigma_{\rm F}$,
between $10^{22}$ and $3{\times}10^{22}$\,cm$^{-2}$, the correction enlarges the error of the uncorrected image by up to $1.4$
times. Where $f$ declines, at the highest surface densities, it reduces the error slightly.

% =====================================================================


\section{Validation}
\label{sec:valid}

A correction can be tested only where the true surface density is known. The correction was tested in three settings, all computed
with radiative transfer. The first is a grid of models of cores, each in its own uniform cloud. The second is a set of isolated
filaments. The third is a simulated sky, in which a cloud, filaments and cores overlap along the lines of sight. All of them were
computed for $G_0 = 10$, so $g = 1$ in Eq.~(\ref{eq:factorg0}). Holding the factor at its maximum changes it only above
$10^{22.6}$\,cm$^{-2}$. Within the surface densities of these models, the change is at most $2$\%, and the results below are given
without it. The correction for other fields was tested separately (Appendix~\ref{sec:field}).

\subsection{On radiative transfer models of cores}
\label{sec:valgrid}

The correction was tested on the $1815$ models of the subdivision grid (Sect.~\ref{sec:modcores}), which were not used to fit the
relation of Eq.~(\ref{eq:law}). The image at $13.5{\arcsec}$ was corrected with the factor evaluated on that image and, for
comparison, with the factor evaluated on the image at $36.3{\arcsec}$. The corrected images were compared with the true surface
density of the model at the same resolution in three ways. Table~\ref{tab:valgrid} gives the results.

The first measure is the accuracy ratio of the total surface density within the footprint of the model, that is, the total in the
corrected image divided by the same total in the true image. The footprint is the area where the model stands above its
surroundings in the true image at $13.5{\arcsec}$. This measure shows whether the correction restores the right amount of material.
The second measure is the mass $M$ of the core. It is the same total with the background subtracted, divided by the same quantity
in the true image. The background is the median surface density in a ring $15{\arcsec}$ wide just outside the footprint. The mass
was measured for $1750$ of the models. For the other $65$ models, the ring contains fewer than $20$ valid pixels, because the
footprint reaches the edge of the image, or the true mass within the footprint is not positive, because the core barely stands
above its cloud. The third measure is the radial profile of the relative residual. The residual is the corrected minus the true
surface density, divided by the true surface density, as a function of the distance from the center of the core. Each measure was
computed for every model. Table~\ref{tab:valgrid} gives the median over the models and the 16th and 84th percentiles, which enclose
$68$\% of the models (one standard deviation from the mean of a Gaussian).

\begin{table}
\caption{Validation on $1815$ radiative transfer models that were not used to fit the relation of Eq.~(\ref{eq:law}). Total: the
accuracy ratio of the surface density summed over the footprint of the model. Core mass: the accuracy ratio of the same sum less
the background, for the $1750$ models in which the mass can be measured. Center and $21{\arcsec}$: the relative residual of the
surface density at $3$ and $21{\arcsec}$ from the center of the core. The rows give the image before the correction and after it,
with the factor evaluated on the image at $36.3{\arcsec}$ or at $13.5{\arcsec}$. The first two columns give the median over the
models, with the 16th and 84th percentiles in parentheses; the last two give medians.}
\label{tab:valgrid}
\centering
\begin{tabular}{lcccc}
\hline\hline
\noalign{\smallskip}
\!\!\! Image & \!\!\!\! Total & \!\!\!\! Core mass & \!\!\!\!\! Center & $21{\arcsec}$\!\!\\
\noalign{\smallskip}
\hline
\noalign{\smallskip}
\!\!\! uncorr.                & \!\!\!\!\! 0.78 (0.56--0.93) & \!\!\!\!\! 0.63 (0.46--0.86) & \!\!\!\!\!\!\! $-0.44$ & \!\!\!\!\!\!\! $-0.21$\!\!\\
\!\!\! $f$ at $36.3{\arcsec}$ & \!\!\!\!\! 0.99 (0.92--1.02) & \!\!\!\!\! 0.94 (0.71--1.09) & \!\!\!\!\!\!\! $-0.15$ & \!\!\!\!\!\!\! $+0.01$\!\!\\
\!\!\! $f$ at $13.5{\arcsec}$ & \!\!\!\!\! 0.99 (0.93--1.02) & \!\!\!\!\! 0.98 (0.76--1.12) & \!\!\!\!\!\!\! $-0.12$ & \!\!\!\!\!\!\! $+0.02$\!\!\\
\noalign{\smallskip}
\hline
\end{tabular}
\end{table}

\begin{figure*}
\centering
\includegraphics[width=0.96\textwidth]{fig_grid_validation.pdf}
\caption{The correction tested on $1815$ radiative transfer models that were not used to fit the relation of Eq.~(\ref{eq:law}).
\emph{Left}: the relative residual of the surface density, the reconstructed minus the true value divided by the true value,
against the distance from the center of the core, as the median over the models. The dashed line marks half the beam of
$13.5{\arcsec}$. \emph{Right}: the accuracy ratio of the core mass, for the $1750$ models in which it can be measured, when the
footprint within which it is measured is too small (negative abscissa) or too large (positive abscissa) by the given fraction of
its equivalent radius. The background is always taken in a ring just outside the footprint used. The lines are medians over the
models, and the band marks the 16th to 84th percentiles for the factor evaluated at $36.3{\arcsec}$.}
\label{fig:valgrid}
\end{figure*}

Without the correction, the reconstruction contains $78$\% of the true surface density and $63$\% of the core mass. The
underestimation is largest at the center of the core, where it reaches $44$\%. The correction restores the total to $99$\% and the
core mass to $98$\% in the median. Beyond $15{\arcsec}$ from the center, the corrected surface density is about $2$\% too high
(Fig.~\ref{fig:valgrid}, left). This is the uniform cloud around the core, which the calibration of Sect.~\ref{sec:diffuse}
corrects slightly too much. An underestimation of $12$\% remains at the center of the core, at $3{\arcsec}$. No correction that
depends on the surface density alone can remove it. The center of a core is colder than the center of a uniform structure with the
same surface density, and the surface density does not reveal this. The factor evaluated at $36.3{\arcsec}$ gives almost the same
totals but a larger underestimation at the center (Sect.~\ref{sec:base}).

The spread over the models is larger than these medians. For one core in six, the corrected mass is more than $24$\% too low. For
another one in six, it is more than $12$\% too high. A correction that depends on the surface density alone cannot remove this
spread. Models with the same surface density can differ in how their temperature varies along the line of sight.

In a real image, the footprint of a core is not known. It must be estimated, together with the background beneath the core. The
right panel of Fig.~\ref{fig:valgrid} shows how an error in the footprint changes the measured mass. A footprint that is too small
by a quarter of its radius lowers the median corrected mass from $0.98$ to $0.84$ of the true mass. A footprint that is too small
by half its radius lowers it to $0.52$, because part of the core is then counted as background. A footprint that is too large by a
quarter or a half of its radius raises the median mass to $1.07$ or $1.13$, and the spread over the models widens sharply. These
errors are lower limits. The models lie in a cloud of uniform density, whose surface density declines smoothly with radius. A real
core is surrounded by filaments and other cores. They enter the background ring in ways that no model reproduces.

\subsection{On radiative transfer models of filaments}
\label{sec:cyl}

The seven cylinders of Sect.~\ref{sec:modfil} were reconstructed in the same way. They are isolated, so their
profiles are not confused by any background. Figure~\ref{fig:cylinders} and Table~\ref{tab:cyl} give the results, measured on
the median profile over the middle part of each cylinder. The true profile is convolved to the beam of $13.5{\arcsec}$.

\begin{table*}
\caption{The isolated radiative transfer cylinders at $13.5{\arcsec}$. $\Sigma_0$: the true surface density on the axis. Crest: the
accuracy ratio of the surface density on the axis, before and after the correction. $\Lambda$: the accuracy ratio of the linear
density within $\pm300{\arcsec}$, before and after. $H$: the half-maximum width of the true, the uncorrected and the corrected
profile. $\alpha$: the slope of $\log\Sigma$ against $\log r$ between $20$ and $120{\arcsec}$ from the crest, for the same three
profiles.}
\label{tab:cyl}
\centering
\begin{tabular}{lcccc}
\hline\hline
\noalign{\smallskip}
$\Sigma_0$ (cm$^{-2}$) & Crest & $\Lambda$ & $H$ (${\arcsec}$) & $\alpha$ \\
\noalign{\smallskip}
\hline
\noalign{\smallskip}
$3{\times}10^{21}$                  & 0.86 $\rightarrow$ 0.93 & 0.96 $\rightarrow$ 1.01 & \!48\, \,\,\,58\,\, 55  & 0.89\,\, 0.88\,\, 0.90 \\
$1{\times}10^{22}$                  & 0.66 $\rightarrow$ 0.83 & 0.87 $\rightarrow$ 0.97 & \!48\, \,\,\,71\,\, 61  & 0.89\,\, 0.81\,\, 0.87 \\
$3{\times}10^{22}$                  & 0.50 $\rightarrow$ 0.81 & 0.75 $\rightarrow$ 0.95 & \!48\, \,\,\,79\,\, 59  & 0.89\,\, 0.70\,\, 0.85 \\
$5{\times}10^{22}$                  & 0.46 $\rightarrow$ 0.83 & 0.68 $\rightarrow$ 0.95 & \!48\, \,\,\,75\,\, 59  & 0.89\,\, 0.67\,\, 0.88 \\
$1{\times}10^{23}$                  & 0.42 $\rightarrow$ 0.80 & 0.60 $\rightarrow$ 0.94 & \!48\, \,\,\,68\,\, 62  & 0.89\,\, 0.70\,\, 0.90 \\
$3{\times}10^{22}$, $p=2$           & 0.50 $\rightarrow$ 0.81 & 0.73 $\rightarrow$ 0.94 & \!47\, \,\,\,72\,\, 56  & 1.07\,\, 0.86\,\, 1.02 \\
$3{\times}10^{22}$, $H=80{\arcsec}$ & 0.53 $\rightarrow$ 0.89 & 0.71 $\rightarrow$ 0.95 &   85\,      122\,\, 91  & 0.80\,\, 0.59\,\, 0.76 \\
\noalign{\smallskip}
\hline
\end{tabular}
\end{table*}

\begin{figure*}
\centering
\includegraphics[width=0.97\textwidth]{fig_cylinders.pdf}
\caption{The isolated radiative transfer cylinders reconstructed by \emph{hires} at $13.5{\arcsec}$ and corrected. \emph{Left}: the
profile across the cylinder with $3{\times}10^{22}$\,cm$^{-2}$ on its axis, true (convolved to the beam), uncorrected and corrected.
\emph{Right}: relative error of the corrected surface density across all seven cylinders, labeled by their surface density on the 
axis. Each profile is the median over the middle part of the cylinder.}
\label{fig:cylinders}
\end{figure*}

Without the correction, the crests of the filaments are $14$ to $58$\% too low. The underestimation grows with the surface density
on the axis. The linear densities are $4$ to $40$\% too low. The profiles are too wide and too shallow, because the crest is
suppressed much more than the wings. After the correction, the linear densities are within $6$\% of the truth, and the slopes of
the profiles are within $0.05$ of the true ones. The crests remain $7$ to $20$\% too low. This is the same limit as at the centers
of the cores: the axis of a filament is colder than the middle of a uniform structure with the same surface density. For the same
reason, the corrected widths $H$ are still $15$ to $30$\% too large for the filaments of $40{\arcsec}$. For the filament of
$80{\arcsec}$, which is well resolved, the width is $7$\% too large. Beyond about $20{\arcsec}$ from the crest, the corrected
profiles are within $7$\% of the truth.

\subsection{On a simulated sky}
\label{sec:sim}

The models of Sects.~\ref{sec:valgrid} and \ref{sec:cyl} each contain a single structure. An observed field does not. The
simulated sky of Sect.~\ref{sec:modcloud} contains a cloud, seven filaments and $96$ cores, which overlap along the lines of sight.

\begin{figure*}
\centering
\includegraphics[width=\textwidth]{fig_simsky6.pdf}
\caption{The simulated sky, built from radiative transfer models of a cloud, seven filaments and $96$ cores. \emph{Upper left}: the
true surface density at $13.5{\arcsec}$. \emph{Upper middle}: the relative error of the corrected surface density, (corrected $-$
true) / true. \emph{Upper right}: the same for a sky that contains the cloud alone. \emph{Lower panels}: the relative error
across each filament, as the median over the outer rows of the field, where the cloud is weak, and over its central rows, where
the cloud is dense. The filaments are labeled by their surface density on the axis. The cores are masked in the profiles. The
dashed curve is the cloud alone at the same positions.}
\label{fig:simsky6}
\end{figure*}

Figure~\ref{fig:simsky6} shows the result for each component. The cloud alone, in a sky without filaments and cores, is corrected
to $0.89$ of the truth in the median. It is $6$\% too low at $10^{21}$\,cm$^{-2}$ and $14$\% too low at $10^{22}$\,cm$^{-2}$.
Before the correction, the reconstruction contains $0.76$ of the truth. The fractional underestimation is the same at
$13.5{\arcsec}$ and at $36.3{\arcsec}$, within $0.5$\%. The
remaining underestimation of $\Sigma$ in the cloud is the price of the average of Sect.~\ref{sec:diffuse}, which gives half the
weight to the sphere.

In the outer parts of the field, the cloud is a few $10^{21}$\,cm$^{-2}$. There the crests of the filaments are corrected to within
$4$ to $15$\%, and their flanks to within $5$\%. Their outer wings are corrected too little, by $6$ to $10$\%, as much as the cloud
around them. In the center of the field, the cloud reaches $2$ to $3\times10^{22}$\,cm$^{-2}$. There the flanks of the dense
filaments are corrected too much, by $10$ to $19$\%, while the cloud alone at the same positions is corrected too little. Both
effects come from the superposition of two structures, which a correction based on the surface density alone treats as one. They
have the signs predicted in Appendix~\ref{app:arrange}: a structure in a weak cloud is under-corrected, and a structure in a dense
cloud is over-corrected. The central cloud of the simulated sky is two to three times denser than the background of the filaments
in Aquila, which \emph{getsf} finds to be $6$ to $10\times10^{21}$\,cm$^{-2}$ at most. In a field like Aquila, the over-correction
of the flanks should therefore be about $10$\%. Over the whole field, with the cores masked, the corrected surface density is
$0.96$ of the truth.

\subsection{Where the correction fails, and the resolution of the factor}
\label{sec:base}

Two kinds of structure are not corrected well by any correction that depends on the surface density alone. The centers of
concentrated structures, the centers of cores and the crests of filaments, remain $7$ to $20$\% too low. Cores narrower than two
beams remain up to $40$\% too low at their centers, surrounded by an over-corrected ring, because the beam spreads their cold
centers over an area of lower surface density. Superposed structures of very different temperatures are corrected too little or too
much, by $10$ to $20$\% (Sect.~\ref{sec:sim}). Corrections that also read the fitted temperature do not help
(Appendix~\ref{app:tcorr}).

The factor can be evaluated at any resolution. Evaluated on the image at $13.5{\arcsec}$, it gives core masses $5$\% higher than
evaluated at $36.3{\arcsec}$, closer to the truth, and a smaller underestimation at the centers (Table~\ref{tab:valgrid}). On the
cloud of the simulated sky, the two agree to $0.5$\%. The factor was therefore evaluated on the image that is corrected. A factor
evaluated at $36.3{\arcsec}$ remains a conservative alternative.

% =====================================================================

\section{Application to the Aquila cloud}
\label{sec:aquila}

The correction was applied to a $31\arcmin \times 31\arcmin$ sub-field of the Aquila cloud (centered at RA $277{\fdg}4765227$, Dec
$-1{\fdg}9962943$). The field was observed with \emph{Herschel} \citep{Andre+2010, Konyves+2015}. The reconstruction used the
images at $160$, $250$, $350$, and $500${\um}. The field contains $625\times 625$ pixels of $3{\arcsec}$. The factor was evaluated
on the reconstruction at $13.5{\arcsec}$ and applied to it. Every pixel was corrected, including those toward protostars
(Appendix~\ref{sec:notdo}).

The radiation field of the cloud was estimated with the procedure of Appendix~\ref{app:fieldest}, from the diffuse material of the
whole Aquila map, which contains $29$ million pixels. The field is $G_0 \approx 33$, with $25$ and $41$ for the 16th and 84th
percentiles of the temperature. The cloud is therefore heated by a field about three times stronger than the field of the
calibration, and the correction used $G_0 = 33$ in Eq.~(\ref{eq:factorg0}).

The median correction factor over the field is $1.68$. The factor is below $1.20$, $1.54$, $1.91$ and $2.28$ for $1$, $25$, $75$
and $99$\% of the pixels, and it reaches its maximum, $2.28$, above $3{\times}10^{22}$\,cm$^{-2}$. The correction is therefore
significant everywhere in the field. The total surface density of the field increases by a factor of $1.89$. With the calibration
field $G_0 = 10$, the median factor would be about $1.45$ and the total would increase by a factor of about $1.62$.

The distribution of surface densities changes more than the median factor suggests, because the factor rises with the surface
density. The median surface density of the field increases from $9.7{\times}10^{21}$ to $1.6{\times}10^{22}$\,cm$^{-2}$, and its
99th percentile from $5.6{\times}10^{22}$ to $1.3{\times}10^{23}$\,cm$^{-2}$. The number of pixels above $10^{22}$\,cm$^{-2}$
increases from $178\,641$ to $318\,215$. Above $3{\times}10^{22}$\,cm$^{-2}$ it increases from $17\,041$ to $85\,343$, and above
$10^{23}$\,cm$^{-2}$ from $666$ to $7049$, a factor of $10.6$. A much larger part of the cloud therefore exceeds any given
threshold after the correction. This matters for any criterion based on a surface density threshold, such as a threshold for star
formation.

The corrected surface density implies a lower dust temperature. It was defined as the temperature at which the corrected surface
density gives the same intensity at $500${\um} as the uncorrected surface density with the fitted temperature. Its median over the
field is $12.2$\,K, against $15.0$\,K for the fitted temperature. It is below $9.3$ and $18.5$\,K for $1$ and $99$\% of the pixels.

\begin{figure*}
\centering
\includegraphics[width=\textwidth]{fig_aquila.pdf}
\caption{The Aquila field at $13.5{\arcsec}$, corrected for the field $G_0 = 33$. \emph{Upper panels}: the surface density before
and after the correction, on a common logarithmic scale. \emph{Lower panels}: the image of the correction factors $\mathcal{F}$ and
the corrected dust temperature (Sect.~\ref{sec:aquila}). The factor is smallest in the diffuse material and largest along the
crests of the filaments and at the cores.}
\label{fig:aquila}
\end{figure*}

Figure~\ref{fig:aquila} shows the field before and after the correction, together with $\mathcal{F}$ and the corrected dust
temperature. The correction is not a uniform rescaling. The factor is smallest in the diffuse material between the structures and
largest along the crests of the filaments and at the cores. The correction therefore raises the contrast of the image as well as
its overall level. Because the factor never decreases with the surface density, the densest crests and cores receive the largest
correction, and the corrected temperature shows no artificial dips toward them. The field is not uniform across the cloud. The warm
corner at the upper left of the temperature map is heated more strongly than the rest of the field, and there the single field of
the cloud underestimates the true one. Appendix~\ref{sec:field} discusses the uncertainty due to the radiation field.

\section{Conclusions}
\label{sec:conclusions}

Surface density images of molecular clouds are derived by fitting one modified blackbody to each pixel. These images are
systematically too low wherever the dust temperature varies along the line of sight. The underestimation exceeds a factor of two in
the densest material.

The underestimation cannot be measured from the observed colors. Its signature in the spectrum is $1$ to $6$\% of the intensities,
while the underestimation itself is $9$ to $38$\%. A method that derived one from the other would amplify the photometric errors
$7$ to $9$ times. The correction must therefore come from radiative transfer.

Radiative transfer provides a single relation between the dust temperature and the effective column $\Sigma_{\rm A}$ that shields
the dust. It describes $2086$ models of embedded Bonnor--Ebert spheres and seven models of filaments, with a scatter of $4.6$\%.
Because $\Sigma_{\rm A}$ is an average over all directions of the incoming radiation, a sphere, a cube and a plane-parallel layer
of the same surface density need factors within $10$\% of each other. The relation gives the correction factor $f$ as an explicit
function of the observed surface density. At low surface densities, the factor is calibrated on a radiative transfer model of an
extended cloud, whose warm surface layers the relation does not describe.

The correction applies to a conventional surface density map at the resolution of the longest waveband and to the
higher-resolution images of the \emph{hires} reconstruction. The
factor was evaluated on the finest image, which gives the sharpest corrected images and the most accurate masses.

The correction was tested on three kinds of radiative transfer models. On $1815$ models of cores that were not used in the
calibration, it raises the median core mass from $0.63$ to $0.98$ of the true mass and the total surface density from $0.78$ to
$0.99$. On seven isolated filaments, it restores the linear densities to $0.94$--$1.01$ of the truth and the slopes of the profiles
to within $0.05$. On a simulated sky of cloud, filaments and cores, all computed with radiative transfer, it restores the total
surface density to $0.96$. The centers of cores and the crests of filaments remain $7$ to $20$\% too low, because they are colder
than a uniform structure with the same surface density. Structures superposed on a cloud are corrected too little where the cloud
is weak and too much where it is dense, by about $10$\% in a cloud like Aquila.

The correction leaves only a few percent in the median core mass, less than the error of a footprint wrong by a quarter of its
radius ($-14$ to $+9$\%). A correction that depends on the surface density alone is therefore sufficient. Corrections that also use
the fitted temperature are not worth their complexity, because a warm pixel may need either more or less correction.

The correction was calibrated for $G_0 = 10$. Models computed for nine fields from $G_0 = 1$ to $483$ show that, without further
adjustment, it is too large by up to $27$\% in the weakest fields and too small by up to $20$\% in fields of a few tens. A field
factor measured on these models, a function of the surface density and the field, reduces these errors to about $10$\%, also for
fields left out of its construction. The field of a cloud is estimated once, from the temperatures of its diffuse material compared
with the models at the same surface density. Nine regions require fields from below $1$ (Taurus) to about $250$ (Cygnus~X). The
other main uncertainty is the floor temperature $T_{\rm min}$ of fully shielded dust, to which $f$ is very sensitive above
$10^{22}$\,cm$^{-2}$.

Applied to a field of the Aquila cloud, whose diffuse material requires $G_0 \approx 33$, the correction raises the total surface
density by a factor of $1.89$ and the number of pixels above $10^{23}$\,cm$^{-2}$ by a factor of $10.6$. The median dust
temperature implied by the corrected surface densities is $12.2$\,K, against $15.0$\,K for the fitted one.

For the adopted dust properties, the total uncertainties of corrected surface densities are about $30$\% in diffuse material and
$50$\% in dense structures, and those of corrected temperatures about $2$\,K (Appendix~\ref{sec:accuracy}). The dust opacity adds a
factor of about $1.5$ to $2$.

\begin{acknowledgements} 
The images of the molecular clouds used in this work were downloaded from the \textit{Herschel} Science
Archive\footnote{\url{http://archives.esac.esa.int/hsa/whsa/}} and processed with the \emph{hires} method
\citep{Menshchikov2021method} by Guo-Yin Zhang. The images were obtained in the \textit{Herschel} Gould Belt Survey
(HGBS\footnote{\url{http://gouldbelt-herschel.cea.fr}}) (PI: Ph.\,Andr{\'e}), HOBYS\footnote{\url{http://hobys-herschel.cea.fr}}
(PI: F.\,Motte, A.\,Zavagno, S.\,Bontemps), and A-CMC (PI: P.\,Harvey) projects. HGBS and HOBYS are the \textit{Herschel} Key
Projects jointly carried out by SPIRE Specialist Astronomy Group 3 (SAG3), scientists of several institutes in the PACS Consortium
(e.g., CEA Saclay, INAF-IAPS Rome, LAM/OAMP Marseille), and scientists of the \textit{Herschel} Science Center (HSC).
\end{acknowledgements}

\bibliographystyle{aa}
\bibliography{hires2_correction}

% =====================================================================

\appendix
\nolinenumbers 

\section{Uncertainties and limitations}
\label{sec:limits}

Table~\ref{tab:accuracy} collects the accuracies measured on the models of this paper, before and after the correction. The
subsections that follow discuss each source of error. The first of them, the background and the footprint of a source, is not an
error of the correction, but it is the largest term for the masses of cores.

\subsection{Summary of the accuracy}
\label{sec:accuracy}

\begin{table*}[ht!]
\caption{Accuracy of surface densities, masses and temperatures derived from \emph{Herschel} images at $160$--$500${\um}, measured
on the radiative transfer models of this paper. Accuracy ratios (measured over true values) and relative errors are medians over
the models, with the 16th and 84th percentiles in parentheses where they are available; ranges span the models listed. For the
diffuse cloud, the range is that between $10^{22}$ and $10^{21}$\,cm$^{-2}$. Uncorrected: the reconstruction of \emph{hires} at
$13.5{\arcsec}$, or a single modified blackbody fit. Corrected: the same after the correction of Eq.~(\ref{eq:correction}). $R_{\rm
eq}$ is the equivalent radius of the footprint of a core.}
\label{tab:accuracy}
\centering
\setlength{\tabcolsep}{4pt}
\begin{tabular}{llll}
\hline\hline
\noalign{\smallskip}
Quantity & Uncorrected & Corrected & Models (Section) \\
\noalign{\smallskip}
\hline
\noalign{\smallskip}
diffuse cloud, accuracy ratio of $\Sigma$               & 0.76                 & 0.89 (0.86--0.94)  & cloud model (\ref{sec:sim}) \\
total $\Sigma$ within core footprints, accuracy ratio   & 0.78 (0.56--0.93)    & 0.99 (0.93--1.02)  & 1815 cores (\ref{sec:valgrid}) \\
core mass from the image, accuracy ratio                & 0.63 (0.46--0.86)    & 0.98 (0.76--1.12)  & 1750 cores (\ref{sec:valgrid}) \\
core center ($3{\arcsec}$), relative error of $\Sigma$  & $-0.44$              & $-0.12$            & 1815 cores (\ref{sec:valgrid}) \\
filament crest, accuracy ratio of $\Sigma$              & 0.42--0.86           & 0.80--0.93         & 7 filaments (\ref{sec:cyl}) \\
filament linear density $\Lambda$, accuracy ratio       & 0.60--0.96           & 0.94--1.01         & 7 filaments (\ref{sec:cyl}) \\
filament width $H$ ($40{\arcsec}$), accuracy ratio      & 1.21--1.65           & 1.15--1.29         & 6 filaments (\ref{sec:cyl}) \\
filament slope $\alpha$, error                          & up to $-0.22$        & within $\pm 0.05$  & 7 filaments (\ref{sec:cyl}) \\
superposed structures, relative error of $\Sigma$       & --                   & $-0.13$ to $+0.23$ & analytic (\ref{app:twocomp}) \\
filament flanks in a dense cloud, relative error        & --                   & $+0.10$ to $+0.19$ & simulated sky (\ref{sec:sim}) \\
inclined layers, random, mean relative error            & --                   & $0$ to $+0.24$     & analytic (\ref{sec:incl}) \\
inclined filaments, random, mean relative error         & --                   & $0$ to $+0.09$     & analytic (\ref{sec:incl}) \\
field $G_0 = 1$ to 483 without $g$, relative error      & --                   & $-0.20$ to $+0.27$ & 9 radiation fields (\ref{sec:field}) \\
field $G_0 = 1$ to 483 with $g$, relative error         & --                   & $-0.13$ to $+0.06$ & 9 radiation fields (\ref{sec:field}) \\
local field wrong by a factor of 2.2, relative error    & --                   & $-0.10$ to $+0.13$ & 9 radiation fields (\ref{sec:field}) \\
floor $T_{\rm min}$ wrong by $\pm 1$\,K, relative error & --                   & $-0.22$ to $+0.51$ & relation (\ref{sec:floor}) \\
core footprint wrong by $\pm R_{\rm eq}/4$, mass        & --                   & $-0.14$ to $+0.09$ & 1750 cores (\ref{sec:bgerr}) \\
core footprint wrong by $\pm R_{\rm eq}/2$, mass        & --                   & $-0.47$ to $+0.15$ & 1750 cores (\ref{sec:bgerr}) \\
calibration error of 10\% at 500{\um}, $\Sigma$         & 0.16--0.18           & 0.16--0.23         & uniform bodies (\ref{sec:needs}) \\
core temperature from fluxes, $T - T_{\rm mass}$ (K)    & 0.0 ($-0.4$, $+0.9$) & --                 & 1815 cores (\ref{sec:accuracy}) \\
core mass from fluxes, accuracy ratio                   & 0.87 (0.68--0.99)    & --                 & 1815 cores (\ref{sec:accuracy}) \\
\noalign{\smallskip}
\hline
\end{tabular}
\end{table*}

The numbers of Table~\ref{tab:accuracy} rest on noiseless models, on the simplest backgrounds, and on the dust and the radiation
field of the models. For observers, they can be combined into a realistic total, adding the terms in quadrature. In diffuse
material, the terms are the bias that the correction leaves ($10$\%), the superposition of structures ($10$\%), a local field that
differs from the field of the cloud by a factor of about two ($10$\%), the calibration uncertainty of \emph{Herschel} of $10$\% per
waveband ($17$\%, Sect.~\ref{sec:needs}), and an error of $1$\,K in the floor temperature $T_{\rm min}$ ($10$\%). In dense
structures, the same terms are $15$ to $20$, $15$, $15$, $20$ and $35$\%. The floor depends on the dust properties, which are not
known in a real cloud, and it is the largest term in dense structures (Appendix~\ref{sec:floor}). The total is about $30$\% in
diffuse material and about $50$\% in dense structures. In addition, the absolute dust opacity is uncertain by a factor of about
$1.5$ to $2$ \citep{Roy+2014}, and it may increase in dense cores \citep{Ossenkopf+Henning1994, Juvela+2015}. The opacity scales
all surface densities together, apart from such an increase, and no correction of the temperature bias can reduce its uncertainty.
Uncorrected images have the calibration and opacity uncertainties as well, but they are also biased: they underestimate $\Sigma$ by
$10$ to $25$\% in diffuse material and by factors of $1.3$ to $2.4$ in dense structures. The corrected temperature follows from the
corrected $\Sigma$ through the intensity at $500${\um} (Sect.~\ref{sec:aquila}), independently of the absolute opacity. By
Eq.~(\ref{eq:sensitivity}), an error of $30$ to $50$\% in $\Sigma$ corresponds to about $2$\,K at $11$ to $17$\,K. The fitted
temperatures, in contrast, are too high by $1$ to $4$\,K in dense structures. These estimates hold for fields within the grid of models and for lines of sight without internal heating. Inclined layers and
filaments can be over-corrected by a further $10$ to $25$\% between $10^{21.5}$ and $10^{22}$\,cm$^{-2}$
(Appendix~\ref{sec:incl}). For the mass of a core, the errors of its footprint and background must be
added (Appendix~\ref{sec:bgerr}).

\subsection{The background and the footprint}
\label{sec:bgerr}

The mass of a core is the surface density within its footprint minus the background beneath it. Neither the footprint nor the
background is known in a real image. Section~\ref{sec:valgrid} measured the effect on models with the simplest possible background.
A footprint that is wrong by a quarter of its radius changes the median mass by $-14$ to $+9$\%. A footprint that is wrong by half
its radius changes it by $-47$ to $+15$\%. For footprints that are too large, the spread over the models widens sharply. In a real
cloud, a core lies on filaments and among other cores. The background beneath it is not a smooth continuation of its surroundings
\citep{Menshchikov2016}, and the errors are larger. They are discussed in Sect.~3.4.7 of \citet{Menshchikov2021method} and in
Sect.~4.4 of \citet{Menshchikov2023}. The correction does not change any of this, because it acts on the image, not on the source.
This is the largest uncertainty of a core mass. It is also the reason why the simplest correction was preferred to the more complex
ones of Appendix~\ref{app:tcorr}. None of them changed the median core mass by more than a few percent, which is much less than
this term.

\subsection{The radiation field}
\label{sec:field}

The factor $f$ was calibrated for $G_0 = 10$. In the nine fields of the grid, from $G_0 = 0.98$ to $483$, the correction without
the field factor $g$ errs by up to $+27$\% in weak fields and by up to $-20$\% in fields of a few tens (Appendix~\ref{app:field}).
With $g$, the errors are within $\pm 10$\% up to $6{\times}10^{22}$\,cm$^{-2}$ and between $-13$ and $+6$\% above. With $g$ built
without the field being tested, they are within $-11$ and $+3$\% (Table~\ref{tab:fieldR}).

The field factor needs the field of the cloud, which is estimated from the diffuse material of the map
(Appendix~\ref{app:fieldest}). The spread of the temperatures of the diffuse material gives the uncertainty of this estimate. In
Aquila, the fields at the 16th and 84th percentiles of the temperature are $25$ and $41$. The field factor changes by $\pm 3$\%
between them. A local field that differs from the field of the cloud by a factor of $2.2$, one step of the grid, changes the
corrected $\Sigma$ by $-10$ to $+13$\%, and a factor of $4.7$ by $-17$ to $+24$\% (Table~\ref{tab:fieldR}, first column). A single
field cannot describe a cloud whose parts are heated very differently, such as the surroundings of OB stars
(Appendix~\ref{app:fieldest}).

\subsection{The arrangement along the line of sight}

The method replaces the structures along a line of sight by a single uniform body with the same total surface density. This is
exact in the plane-parallel limit (Eq.~\ref{eq:invariance}), so the error comes only from the radiation that escapes across the
line of sight (Appendix~\ref{app:arrange}). A concentrated structure needs $2$ to $4$\% more correction than a uniform one for a
central-to-edge density ratio of $10$, and $8$ to $19$\% more for a ratio of $10^3$. A structure embedded in a cloud is corrected
between $13$\% too little, in a weak cloud, and $23$\% too much, near a face of a dense cloud. The depth of a structure within a
cloud cannot be observed, so this uncertainty cannot be reduced by a better calculation.

\subsection{Inclination}
\label{sec:incl}

The factor was computed for lines of sight normal to a layer or perpendicular to the axis of a filament. Seen at an angle $i$ from
that direction, a layer or a filament has the same temperature profile along the line of sight, but its surface density and its
fitted surface density are both larger by the factor $1/\cos i$. The factor that it requires does not change, whereas the
correction applies the factor of a denser structure. With the adopted factor, the corrected $\Sigma$ is too high by $2$, $5$, $13$,
and $27$\% at $i = 30$, $45$, $60$, and $70\degr$ for $\Sigma_{\rm F} = 10^{21.5}$\,cm$^{-2}$, and by $5$, $13$, $23$, and $29$\% for
$10^{22}$\,cm$^{-2}$. Below $10^{21}$\,cm$^{-2}$ the error is a few percent at most, and above $10^{22.5}$\,cm$^{-2}$ it vanishes,
because the factor is held at its maximum there. For random orientations, the mean error is up to $24$\% for layers and up to $9$\%
for filaments, in both cases between $10^{21.5}$ and $10^{22}$\,cm$^{-2}$, where the factor rises most steeply.

\subsection{Diffuse material}

Extended diffuse material is corrected too little, by about $6$\% at $10^{21}$\,cm$^{-2}$ and by up to $14$\% at
$10^{22}$\,cm$^{-2}$, while the uniform clouds around the model cores are corrected about $2$\% too much (Sect.~\ref{sec:diffuse}).
A calibration on the extended cloud alone would over-correct bounded clouds by $5$ to $12$\%, so the average was preferred, which
divides the uncertainty between the two limits.

\subsection{The temperature relation and its floor}
\label{sec:floor}

The size of the correction depends on Eq.~(\ref{eq:law}), and in particular on the temperature $T_{\rm min} = 7.0$\,K that dust
approaches when it is completely shielded. The factor is very sensitive to $T_{\rm min}$. Lowering $T_{\rm min}$ by $1$\,K raises
the factor by $30$\% at $10^{22}$\,cm$^{-2}$ and by $51$\% at $3{\times}10^{22}$\,cm$^{-2}$. Raising $T_{\rm min}$ by $1$\,K lowers
the factor by $13$ and $22$\%, respectively. The floor is therefore the part of the relation that most needs to be right. The fit
to $2086$ spheres and seven cylinders constrains it well (Fig.~\ref{fig:shielding}). Above $10^{23}$\,cm$^{-2}$, the models contain
too little mass to constrain the relation. Other radiative transfer models of prestellar cores also give central temperatures of
about $7$ to $8$\,K in the local radiation field \citep{Evans+2001, Zucconi+2001}. A floor of $10$\,K, as is sometimes assumed,
would substantially lower the factor at the highest surface densities. The geometry adds at most $5$\% (Sect.~\ref{sec:orient}).

\subsection{What the correction does not do}
\label{sec:notdo}

The correction assumes that the dust is heated from outside, as in the models. A line of sight through a protostar contains an
internal source, and its temperature gradient is reversed. The models do not describe such a line of sight, and the accuracy of the
correction there is unknown. Every pixel was nevertheless corrected in the same way. Leaving the protostars uncorrected would
require detecting them first, and the correction would then depend on a source extraction. The masses of protostellar cores
measured in corrected images should therefore be used with caution.

\section{The grids of models}
\label{sec:grid}

All models of cores were Bonnor--Ebert spheres embedded in a uniform spherical cloud, heated from outside by the interstellar
radiation field and computed with \emph{radmc-3d} for the dust of Sect.~\ref{sec:models}. Three grids were used, which are
summarized in Table~\ref{tab:grids}.

The main grid varied three parameters: the surface density $\Sigma_{\rm emb}$ of the embedding cloud, from $3.75\times10^{20}$ to
$9.6\times10^{22}$\,cm$^{-2}$ in $9$ steps of a factor of $2$; the mass of the Bonnor--Ebert sphere, from $1.6{\times}10^{-3}$ to
$31\,M_{\odot}$ in $26$ steps of $0.17$\,dex; and the half-maximum width of its volume density profile, from $0.003$ to $0.25$\,pc
in $19$ steps of $0.11$\,dex. Only the models whose peak surface density exceeded that of the embedding cloud by at least $5$\%
were computed, and one of them, the most extended one, was not completed because it required too much memory. The central-to-edge
density ratios $\rho_{\rm C}/\rho_{\rm e}$ of the $2086$ computed models ranged from $1.04$ to $2.1{\times}10^{3}$. Three more
models extended the grid to $\Sigma_{\rm emb}$ of $1.9$ to $7.9\times10^{23}$\,cm$^{-2}$. They were optically thick in the far
infrared and were not used in the calibration (Sect.~\ref{sec:modcores}). The temperature of each model, tabulated shell by shell,
entered the relation of Eq.~(\ref{eq:law}).

The models of the subdivision grid lay between the models of the main grid and were not used in the relation. The footprint of a
model is the area of its convolved true image that stands above the embedding cloud.

The grid of radiation fields repeated $305$ models of the main grid, together with $34$ models at a tenth embedding surface density
of $1.92\times10^{23}$\,cm$^{-2}$. They were computed for nine strengths of the radiation field, from $G_0 = 0.98$ to $483$, in
steps of a factor of $2.17$. The lowest field is the field of \citet{Black1994} alone, without the hot component, and no weaker
field is possible in this implementation. About five models per field were not computed, because they required the largest images.
They do not affect the results of Appendix~\ref{app:field}.

\begin{table}[ht!]
\caption{The grids of radiative transfer models.}
\label{tab:grids}
\centering
\begin{tabular}{lccl}
\hline\hline
\noalign{\smallskip}
Grid &\!\!\! $G_0$ &\!\!\! Models & Used for\\
\noalign{\smallskip}
\hline
\noalign{\smallskip}
main             &\!\!\! $10$         &\!\!\! $2089$         & relation of Eq.~(\ref{eq:law})\\
subdivision      &\!\!\! $10$         &\!\!\! $1815$         & validation, Sect.~\ref{sec:valgrid}\\
radiation fields &\!\!\! $0.98{-}483$ &\!\!\! $9{\times}339$ & Appendix~\ref{app:field}\\
\noalign{\smallskip}
\hline
\end{tabular}
\end{table}

\section{Radiative transfer models of filaments and of the cloud}
\label{app:rtmodels}

\subsection{Cylinders}

The cylinders were computed with \emph{radmc-3d} on an axisymmetric spherical grid of $400$ cells in radius and $480$ cells in the
polar angle, which covered both sides of the equatorial plane. The grid was refined toward the axis and toward the equatorial
plane. The temperatures of the Monte Carlo calculation were averaged over the middle part of the cylinder, in bins of the distance
$r$ from the axis, which removed most of the Monte Carlo noise. The isolated cylinders were computed with $2\times10^9$ photon
packets.

The mirror symmetry of \emph{radmc-3d} was not used. On a test sphere, it made the temperatures $0.6$ to $1.1$\,K too low, whereas
the full grid agreed with a one-dimensional calculation to $0.15$\,K.

\subsection{Filaments embedded in the cloud}

The filaments of the simulated sky were the cylinders of Sect.~\ref{sec:modfil}, each computed again inside a cylindrical envelope
that stood for the cloud around it. The envelope was uniform in radius between $R = 0.5$\,pc and $1$\,pc, and the cylinder was
$10.3$\,pc long, longer than the field, so that its ends lay outside the images. The surface density of the envelope through the
axis followed the cloud of the sky along the filament, row by row. A cylindrical envelope shields the axis less than a cloud layer
of the same surface density, because the layer also blocks the directions close to its plane. The envelope was therefore made $1.3$
to $1.6$ times denser than the cloud, so that the mean intensity at the axis was the same as in the middle of the cloud layer. The
temperatures were averaged in a running window $0.2$\,pc long along the axis, which followed the variation of the envelope. Each
model was imaged twice: once in full, and once with the volume of the filament emptied while all temperatures were kept. The
difference is the emission of the filament inside the cloud, without the emission of the envelope itself. It was added to the cloud
of the sky, which was assumed to extend far in front of and behind the filament. The embedded filaments were computed with $10^9$
photon packets. Runs with $8\times10^9$ packets gave the same temperatures to $0.02$\,K and the same emission to $0.3$\%.

\subsection{The cloud}

The cloud of the simulated sky was a three-dimensional radiative transfer model whose projection is the map of the cloud. A surface
density map does not determine the density along the line of sight, and the model is one of the many distributions that project to
the map. The map had scale-free fluctuations, with a power spectrum $P(k) \propto k^{-3.0}$, a median surface density of about
$3{\times}10^{21}$\,cm$^{-2}$, and no large-scale gradient. Real clouds such as Aquila have an extended background that rises
toward their centers. The map was therefore multiplied by $1 + 5\,[1 + (r/600{\arcsec})^2]^{-0.75}$, where $r$ is the distance from
the center of the field. This factor is $6$ at the center, nearly constant within $600{\arcsec}$, and declines as $r^{-1.5}$
beyond, toward unity far from the center, while the fine structure of the map is kept. The same enhanced map was used for the cloud
of the simulated sky and for the envelopes of the embedded filaments. The map was averaged into cells of $12{\arcsec}$ and extended
on every side by margins of $2$\,pc, by mirror reflection. The margins kept the heating through the sides of the box away from the
field.

The volume density was built in three steps. First, a Gaussian random field $G$ of unit variance, with the same power spectrum
$P(k) \propto k^{-3.0}$ in three dimensions, filled a box with the area of the extended map and a depth of $2$\,pc. Second, its
log-normal transform $g = \exp(s\,G - s^2/2)$, with the amplitude $s = 0.5$, gave positive fluctuations with a mean of unity and a
root-mean-square contrast of about $50$\%. Third, every column of $g$ was rescaled so that its integral equaled the surface density
of the map in that pixel, $n(x, y, z) = \Sigma(x, y)\,g(x, y, z) / \int g\,{\rm d}z$. The projection of the box is therefore
exactly the map, while the density varies along every line of sight as it does across the sky. With the margins, the box had $712
\times 712 \times 132$ cells ($10.8 \times 10.8 \times 2.0$\,pc) and was computed with $8\times10^{10}$ photon packets. Fewer
photon packets, cells of $24{\arcsec}$, or a box without fluctuations along the line of sight gave the same correction factors to
$0.1$\%.

Figure~\ref{fig:cube3d} gives a view of the model. Figure~\ref{fig:cube} shows a slice through it, along one row of the field and
the line of sight, and compares it with the analytic slab of the same surface density. The surface layers are up to $3$\,K warmer
than the slab, because they are heated by the unattenuated ultraviolet radiation from half the sky, which the gray relation of
Eq.~(\ref{eq:law}) does not describe. Thin columns between dense regions are up to $2$\,K colder than the slab, because their dense
neighbors shield them from the side.

\begin{figure}[ht!]
\centering
\includegraphics[width=\columnwidth]{fig_cube3d.pdf}
\caption{A cut-away view of the radiative transfer model of the cloud, seen from the observer's side. The front face is the surface 
density map of the lower half of the field, oriented as in Fig.~\ref{fig:simsky6}; the depth runs along the line of sight. The 
upper face is a cut through the middle row of the field, and the right face a cut along its last column. Both show the dust 
temperature in the plane of the cut.}
\label{fig:cube3d}
\end{figure}

\begin{figure*}[ht!]
\centering
\includegraphics[width=0.9\textwidth]{fig_cube.pdf}
\caption{A slice through the radiative transfer model of the cloud, along the middle row of the field (vertical axis) and the line
of sight (horizontal axis). The depth along the line of sight is horizontal in all four panels. \emph{Left three panels}: the H$_2$
number density, the dust temperature, and the difference between the dust temperature and that of the analytic slab with the same
surface density in each column. \emph{Right panel}: the dust temperature against the depth for three columns of different surface
density, in the radiative transfer model (solid) and in the analytic slab (dashed).}
\label{fig:cube}
\end{figure*}

As a result, the emission of the radiative transfer model at $160${\um} is up to $1.5$ times that of the analytic slab with the same
surface density. In the simulated sky, the emission of the analytic slab was multiplied in each band by the ratio of the emission of
the radiative transfer model to that of the slab, computed in each cell and interpolated to the pixels. The sky thus had the
temperatures of the radiative transfer model and the fine structure of the map of the cloud.

\section{The arrangement of the material along the line of sight}
\label{app:arrange}

The correction is applied pixel by pixel. A pixel receives the emission of everything along its line of sight, and a single
modified blackbody is fitted to this sum. The fit is therefore biased by the whole distribution of dust temperatures along the line
of sight, not by the temperature of any one structure. The method replaces this distribution by the one of a single uniform body
with the same total surface density. This appendix measures the error of this replacement. In a plane-parallel geometry, the
emergent emission depends only on the total surface density (Eq.~\ref{eq:invariance}).

\subsection{Concentration}
\label{app:concentr}

In the plane-parallel limit, the density profile has no effect at all. Concentration can therefore act only through the radiation
that escapes across the line of sight. A bounded structure allows this escape, and a layer does not. The effect was measured on a
sphere with the density $\rho(r) \propto (1 + r^2/a^2)^{-1}$ out to its boundary. The sphere was normalized so that the column
through its center equaled the observed one, and $a$ was set by the central-to-edge density ratio. The parcels along the central
line of sight were weighted by their density, as in the emission. For each parcel, the column was integrated along the ray to the
boundary in every direction of the average.

\begin{table}[ht!]
\caption{Correction factor required by the central line of sight of a sphere of the given central surface density, for four
central-to-edge density ratios. The first column is a uniform sphere.}
\label{tab:concentr}
\centering
\begin{tabular}{lcccc}
\hline\hline
\noalign{\smallskip}
$\Sigma$ (cm$^{-2}$) & uniform & $10$ & $10^2$ & $10^3$ \\
\noalign{\smallskip}
\hline
\noalign{\smallskip}
$5{\times}10^{21}$ & $1.087$ & $1.124$ & $1.155$ & $1.172$ \\
$1{\times}10^{22}$ & $1.270$ & $1.323$ & $1.392$ & $1.437$ \\
$3{\times}10^{22}$ & $1.722$ & $1.776$ & $1.906$ & $2.045$ \\
$1{\times}10^{23}$ & $1.898$ & $1.940$ & $2.043$ & $2.228$ \\
\noalign{\smallskip}
\hline
\end{tabular}
\end{table}

Table~\ref{tab:concentr} gives the result. Concentration always raises the required correction, because the mass of a concentrated
structure lies where the shielding is deepest. The effect is $2$ to $4$\% for a density ratio of $10$ and $8$ to $19$\% for a ratio
of $10^3$. It is one-sided. A correction built on uniform structures therefore under-corrects a concentrated structure and never
over-corrects it.

\subsection{Two superposed components}
\label{app:twocomp}

The second case is a line of sight that crosses a cloud with a denser structure embedded in it at an unknown depth. The model
consisted of a plane-parallel cloud of surface density $\Sigma_{\rm cl}$. The material at the column depth $x$ in the cloud is
shielded by the angular average of $x/|\mu|$ toward the observer and of $(\Sigma_{\rm cl} - x)/|\mu|$ away from the observer. A
uniform sphere with the surface density $\Sigma_{\rm s}$ through its center is embedded in the cloud. A fraction $q$ of the cloud
lies in front of the sphere. A parcel at the signed position $z$ along the diameter of the sphere, in units of its radius, is
shielded by its own sphere over the path
\begin{equation}
   l(z,\psi) = -z\cos\psi + \left(1 - z^2\sin^2\psi\right)^{1/2},
\label{eq:lsphere}
\end{equation}
in units of the radius. It is also shielded by the cloud over $q\,\Sigma_{\rm cl}/|\cos\psi|$ toward the observer and
$(1-q)\,\Sigma_{\rm cl}/|\cos\psi|$ away from the observer. The angular average of Eq.~(\ref{eq:seff}) was taken over $\psi$ with
the isotropic weight $\sin\psi$, and the temperature of every parcel followed from Eq.~(\ref{eq:law}). The emissions of the cloud
and of the sphere were added with their masses, and one modified blackbody was fitted to the sum, exactly as in the reconstruction.
This gave the fitted surface density $\Sigma_{\rm F}$ of the composite line of sight and the factor $\Sigma/\Sigma_{\rm F}$ that it
requires. The method sees only $\Sigma_{\rm F}$ and applies $f(\Sigma_{\rm F})$ of Eq.~(\ref{eq:factor}). The error of the
corrected surface density is therefore $f(\Sigma_{\rm F})\,\Sigma_{\rm F}/\Sigma - 1$.

\begin{figure*}[ht!]
\centering
\includegraphics[width=0.95\textwidth]{fig_superposition.pdf}
\caption{Superposition of a cloud and a sphere. \emph{Upper left}: The dust temperature along a line of sight of
$2{\times}10^{22}$\,cm$^{-2}$. One curve is for the uniform body that the method assumes. The other curve is for the same surface
density divided equally between a cloud and a uniform sphere at the near side of the cloud ($q = 0$). \emph{Upper right}: The
emission of the two lines of sight, with the surface density $\Sigma_{\rm F}$ that a single modified blackbody fit returns. The two
lines of sight have the same true surface density. \emph{Lower left}: The factor that the method applies and the range of factors
that the composite lines of sight require, both against the observed (fitted) surface density. The range covers cloud surface
densities $\Sigma_{\rm cl}$ of $10^{21}$, $3{\times}10^{21}$ and $10^{22}$\,cm$^{-2}$ and all positions of the sphere. \emph{Lower
right}: Relative error of the corrected surface density against the fraction $q$ of the cloud in front of the sphere. Colors give
the surface density of the cloud $\Sigma_{\rm cl}$; solid and dashed curves are for total surface densities of $2{\times}10^{22}$
and $5{\times}10^{22}$\,cm$^{-2}$.}
\label{fig:superposition}
\end{figure*}

Figure~\ref{fig:superposition} gives the result. The error of the corrected surface density ranges from $-13$ to $+23$\%. Its sign
depends on the cloud. A weak cloud ($\Sigma_{\rm cl} = 10^{21}$\,cm$^{-2}$) leaves the sphere nearly isolated. The sphere needs more
correction than the average of a sphere and a layer, so the correction is too small by $8$ to $13$\%, at any depth. A dense cloud
($\Sigma_{\rm cl} = 10^{22}$\,cm$^{-2}$) with the sphere at one of its sides makes the line of sight warmer than a single body. The
correction is then too large, by up to $23$\%. When the sphere lies near the middle of a dense cloud, the error is only a few
percent. The simulated sky of Sect.~\ref{sec:sim}, built entirely from radiative
transfer models, shows both regimes: the wings of filaments in its weak outer cloud are corrected too little, and the flanks of
filaments in its dense central cloud too much.

The upper panels of Fig.~\ref{fig:superposition} show why the effect exists. A sphere at the near side of the cloud loses radiation
sideways, and nothing in front of it blocks the incoming radiation. This line of sight is warmer and brighter than the uniform body
of the same surface density. A single modified blackbody fit returns $1.47{\times}10^{22}$\,cm$^{-2}$ for it and
$1.26{\times}10^{22}$\,cm$^{-2}$ for the uniform body. The composite line of sight requires a factor of $1.36$, but the method sees
$1.47{\times}10^{22}$\,cm$^{-2}$ and applies the factor that a uniform body with this observed surface density requires, which is
$1.67$. The corrected surface density is then $23$\% too high. Nothing observable distinguishes the two lines of sight, so this
error is an irreducible uncertainty of any method of this kind.

\section{Why the observed colors cannot supply the correction}
\label{app:colors}
\label{sec:amplification}

One could try to measure the temperature structure directly, by fitting more than one component to the observed spectrum of each
line of sight. Without noise, this works. A free two-temperature fit from $160$ to $500${\um} recovers $0.96$ to $1.00$ of the true
surface density at the centers of the four models, against $0.62$ to $0.91$ for the single-temperature fit. With noise, it does not
work. Table~\ref{tab:signature} shows why.

\begin{table}[ht!]
\caption{Signature of the temperature mixing in the four-band spectral energy distribution at the center of each model, compared
with the underestimation of $\Sigma$ it produces. Columns two to five give the relative residuals, in percent, of the best
single-component fit to the noiseless intensities; column six their root-mean-square; column seven the fraction of the true surface
density that the single-component fit recovers.}
\label{tab:signature}
\centering
\begin{tabular}{lcccccc}
\hline\hline
\noalign{\smallskip}
$\rho_{\rm C}/\rho_{\rm e}$ &\!\!\!\! 160\,$\mu$m &\!\!\!\! 250\,$\mu$m &\!\!\!\! 350\,$\mu$m &\!\!\!\! 500\,$\mu$m &\!\!\!\! rms &\!\!\! recovery\\
\noalign{\smallskip}
\hline
\noalign{\smallskip}
1.18  &\!\!\!\! $+0.92$ &\!\!\!\! $-1.37$ &\!\!\!\! $-0.44$ &\!\!\!\! $+0.93$ &\!\!\!\! 0.97 &\!\!\!\ 0.914\\
1.28  &\!\!\!\! $+0.88$ &\!\!\!\! $-2.16$ &\!\!\!\! $-0.46$ &\!\!\!\! $+1.82$ &\!\!\!\! 1.50 &\!\!\!\ 0.869\\
4.47  &\!\!\!\! $+3.46$ &\!\!\!\! $-6.44$ &\!\!\!\! $-1.56$ &\!\!\!\! $+5.39$ &\!\!\!\! 4.61 &\!\!\!\ 0.680\\
90.8  &\!\!\!\! $+3.76$ &\!\!\!\! $-8.04$ &\!\!\!\! $-1.64$ &\!\!\!\! $+7.27$ &\!\!\!\! 5.79 &\!\!\!\ 0.619\\
\noalign{\smallskip}
\hline
\end{tabular}
\end{table}

The residuals have the same pattern for every model: positive at $160$ and $500${\um} and negative at $250$ and $350${\um}. This is
the broadening of the spectrum by a mixture of temperatures. The underestimation divided by the root-mean-square residual is $8.9$,
$8.7$, $6.9$ and $6.6$. Any method that infers the underestimation from the colors therefore amplifies the relative photometric
errors by a factor of $7$ to $9$. The calibration uncertainties of \emph{Herschel} observations of Aquila are $20$\% at $70$ to
$160${\um} and $10$\% at $250$ to $500${\um} \citep{Konyves+2015}. For the two weakly concentrated models, the whole signature is
below the photometric accuracy, and a test that accepts a second component only where it significantly improves the fit never
accepts it (Fig.~\ref{fig:mixing}). The correction must therefore be calibrated on models in which the answer is known.

\begin{figure*}[ht!]
\centering
\includegraphics[width=0.96\textwidth]{fig_mixing.pdf}
\caption{\emph{Left}: the emission of a line of sight of $3{\times}10^{22}$\,cm$^{-2}$, the sum over depth of material at every
temperature, with the single modified blackbody fitted to it in the four wavebands and with the weights that the reconstruction
uses. The fit is drawn toward the warmer and brighter material, so the temperature it returns is too high and the surface density
that follows from it is too low. \emph{Right}: the relative residuals of that fit, waveband by waveband, for three total surface
densities, with the fraction of the true surface density recovered beside each. The pattern is always the same, positive at the
ends and negative in the middle, and it grows with the surface density, but it never rises above the calibration uncertainties of
$20${\%} at $160${\um} and $10${\%} at the three longer wavebands.}
\label{fig:mixing}
\end{figure*}

\section{The radiation field}
\label{app:field}

The relation of Eq.~(\ref{eq:law}) was measured for a single field, $G_0 = 10$. A stronger field heats the dust at every
attenuating column, and this changes the underestimation of $\Sigma$ by the fit. To measure the effect, the correction was tested
on the grid of radiation fields of Appendix~\ref{sec:grid}, comparing for every pixel the factor that restores its true surface
density with the factor that the correction applies. The first column of Table~\ref{tab:fieldR} gives the relative error of the
corrected surface density without the field factor. The error changes sign twice: the correction is too large in the weakest
fields, too small in fields of a few tens, and again too large at the highest surface densities in the strongest fields. The
underestimation of $\Sigma$ by a single-temperature fit depends on how the temperature varies along the line of sight, and a
stronger field changes this variation differently at different depths. No simple rescaling of Eq.~(\ref{eq:law}) with the field can
therefore reproduce the errors.

\subsection{The field factor $g$}
\label{app:fieldR}

The field factor $g(\Sigma_{\rm F}, G_0)$ of Eq.~(\ref{eq:factorg0}) was measured on the same grid. For each field, the factor that
each pixel of each model required was tabulated against the fitted surface density. Its logarithm was averaged over the pixels in
bins of $0.1$\,dex of $\Sigma_{\rm F}$ and smoothed over three bins. The field factor is the ratio of this average to the same
average at $G_0 = 10$. Bins with fewer than $2{\times}10^4$ pixels took the value of the nearest populated bin, and above
$10^{22.6}$\,cm$^{-2}$ the field factor was held at its value there. Table~\ref{tab:fieldR} gives the errors of the corrected
surface density, without and with the field factor, and with the field factor built without the field being tested.

\begin{table}[ht!]
\caption{Relative error of the corrected surface density, in percent, for the models of the grid of radiation fields, in four
ranges of $\log_{10}\Sigma_{\rm F}$ (cm$^{-2}$): $21.6$--$22.0$, $22.0$--$22.4$, $22.4$--$22.8$ and $22.8$--$23.4$. Each entry is
the median over the pixels in the range. The rows give the correction without the field factor $g$, with $g$ built from all nine
fields, and with $g$ built without the field tested (left out). The factor was held at its maximum in all rows.}
\label{tab:fieldR}
\centering
\setlength{\tabcolsep}{4pt}
\begin{tabular}{rccc}
\hline\hline
\noalign{\smallskip}
$G_0$ & Without $g$ & With $g$ & Left out \\
\noalign{\smallskip}
\hline
\noalign{\smallskip}
0.98             & $+15$ $+27$ $+21$ $+16$ & $+4$ $-2$ $-5$ $-9$  & -- \\
2.12             & $+14$ $+22$ $+15$ $+10$ & $+2$ $-3$ $-4$ $-9$  & $+2$ $-1$ $-3$ $-8$ \\
4.61             & $+8$ $+11$ $+7$ $-2$    & $+3$ $-1$ $-4$ $-10$ & $+3$ $-1$ $-3$ $-11$ \\
10.0\phantom{0}  & $+1$ $-2$ $-5$ $-13$    & $+1$ $-2$ $-5$ $-13$ & -- \\
21.7\phantom{0}  & $-7$ $-12$ $-12$ $-19$  & $+1$ $-3$ $+3$ $-5$  & $0$ $-5$ $-4$ $-11$ \\
47.1\phantom{0}  & $-12$ $-19$ $-19$ $-20$ & $+1$ $-4$ $-2$ $-3$  & $-1$ $-9$ $-7$ $-8$ \\
102.4\phantom{0} & $-12$ $-15$ $-13$ $-12$ & $+1$ $-4$ $-3$ $-2$  & $-2$ $-7$ $-5$ $-4$ \\
222.2\phantom{0} & $-8$ $-7$ $+1$ $+5$     & $0$ $-3$ $-1$ $+3$   & $0$ $-6$ $-2$ $+1$ \\
482.5\phantom{0} & $-1$ $+6$ $+17$ $+24$   & $0$ $-3$ $0$ $+6$    & -- \\
\noalign{\smallskip}
\hline
\end{tabular}
\end{table}

With the field factor, the error no longer depends on the field. In every field, it is close to the error at $G_0 = 10$, and within
$\pm 10$\% except in the densest range. The remaining underestimation of up to $13$\% in the densest range is that at the centers
of concentrated cores. When a field is left out, its errors grow by a few percent, which measures the error of interpolating
between the fields of the grid.

\subsection{Estimating the field that a cloud requires}
\label{app:fieldest}

The field of a cloud is estimated from the temperatures that the reconstruction fits, by comparing them with those of the models at
the same surface density. For each field of the grid, the median fitted temperature of the models is tabulated in bins of
$0.2$\,dex of the fitted surface density. At a fixed surface density, this temperature rises monotonically with the field. An
observed temperature is therefore converted into the field that the models need to reach it, by interpolating $\log G_0$ against
the median temperature. The procedure uses only the surface density and temperature images at the resolution of the harvested
models, $13.5{\arcsec}$, and it is carried out by the script \texttt{g0\_requirement.py}.

The simple scaling $T \propto G_0^{1/6}$ \citep{IvezicElitzur1997, Tielens2005} for the adopted dust with $\kappa_{\lambda}\propto
\lambda^{-2}$, is not used. It holds for dust that is optically thin to the heating radiation. At a fixed surface density, the
fitted temperature of the models rises only as $G_0^{0.09}$ to $G_0^{0.13}$, because the shielded dust is heated by an attenuated
and redder field and partly by the floor component. The simple scaling would place a cloud at a field several times too weak.

The field is taken from the diffuse material only, below $10^{21.8}$\,cm$^{-2}$. In denser bins, the fields that the procedure
gives rise steadily with the surface density in every cloud. Table~\ref{tab:g0aquila} shows this for Aquila, where the median field
rises from about $20$--$40$ in the diffuse bins to more than $100$ in the densest ones. The rise is found even in Scorpius, which
contains no strong heating source. It is therefore a bias, not a real increase of the field. A dense structure seen through warm
diffuse material has a higher fitted temperature than an isolated core of the same surface density, and the procedure attributes
the extra warmth to the field. Pixels warmer than the strongest field of the grid allows are counted separately. They are heated
internally, by protostars, or by nearby massive stars beyond the range of the grid.

\begin{table}[ht!]
\caption{The procedure of Appendix~\ref{app:fieldest} applied to the whole Aquila map at $13.5{\arcsec}$. Pixels: the number of
pixels in each bin of $\log_{10}\Sigma$, with $\Sigma$ in cm$^{-2}$. $T$: the fitted temperature at the 16th, 50th and 84th
percentiles of the pixels in the bin. $G_0$: the fields that reproduce these temperatures; $^{>}482$ marks a temperature above the
strongest field of the grid. Last column shows the percentage of pixels above the strongest field.}
\label{tab:g0aquila}
\centering
%\footnotesize
\setlength{\tabcolsep}{2.2pt}
\begin{tabular}{lrccl}
\hline\hline
\noalign{\smallskip}
$\log_{10}\Sigma$ & Pixels & $T$ (K) & $G_0$ &\quad\%\\
\noalign{\smallskip}
\hline
\noalign{\smallskip}
21.0--21.2 & 370582   &\, 19.2\,\,19.7\,\,20.0 &\, 21\,\,25\phantom{0}\,\,27     &\phantom{00}0.0\\
21.2--21.4 & 9527486  &\, 18.4\,\,19.0\,\,19.4 &\, 16\,\,21\phantom{0}\,\,24     &\phantom{00}0.0\\
21.4--21.6 & 14764152 &\, 17.7\,\,18.4\,\,19.1 &\, 29\,\,38\phantom{0}\,\,50     &\phantom{00}0.0\\
21.6--21.8 & 3573950  &\, 16.7\,\,17.4\,\,17.9 &\, 47\,\,63\phantom{0}\,\,78     &\phantom{00}0.6\\
21.8--22.0 & 837305   &\, 15.7\,\,16.6\,\,18.3 &\, 56\,\,87\,\,205               &\phantom{00}8.2\\
22.0--22.2 & 242512   &\, 14.6\,\,15.5\,\,17.8 &\, 51\,\,81\,\,299               &\phantom{0}10.9\\
22.2--22.4 & 76237    &\, 13.8\,\,14.5\,\,16.5 &\, 62\,\,94\,\,284               &\phantom{0}10.2\\
22.4--22.6 & 23438    &\, 13.3\,\,13.9\,\,16.0 &\, 74\,\,94\,\,295               &\phantom{0}11.3\\
22.6--22.8 & 8368     &\, 12.8\,\,13.5\,\,15.8 &\, 73\,\,118\,\,$^{>}$482        &\phantom{0}17.6\\
22.8--23.0 & 2874     &\, 12.0\,\,13.5\,\,15.5 &\, 109\,\,252\,\,$^{>}$482\,\,\, &\phantom{0}28.0\\
23.0--23.2 & 546      &\, 13.6\,\,14.6\,\,15.9 & 343\,\,$^{>}$482\,\,$^{>}$482   &\phantom{0}69.8\\
\noalign{\smallskip}
\hline
\end{tabular}
\end{table}

The field of a cloud is the median, weighted by the number of pixels, of the fields of its diffuse bins. The fields at the 16th and
84th percentiles of the temperature give its spread. Table~\ref{tab:g0regions} lists the result for nine regions. Their diffuse
fields span more than two orders of magnitude, from below $1$ in Taurus and about $5$ in California and Perseus to about $250$ in
Cygnus~X. Only Scorpius and Orion~A are close to the field $G_0 = 10$ of the calibration. In Taurus, the diffuse dust ($13.6$ to
$14.4$\,K between $10^{21}$ and $10^{21.4}$\,cm$^{-2}$) is colder than the models in the weakest field of the grid ($14.9$\,K), and
the procedure gives only an upper limit. The correction then uses the field factor of $G_0 = 0.98$. Between $G_0 = 0.98$ and
$2.12$, the errors of the correction without $g$ change by only $1$ to $6$\% (Table~\ref{tab:fieldR}), so this approximation
changes the corrected $\Sigma$ by a few percent at most. Taurus is heated by about the local interstellar field, of order unity in
Habing units, so the estimated fields are not physical field strengths. They would all change if the real dust emitted more or less
efficiently than the dust of the models. The correction is consistent nevertheless, because $g$ was measured on the same dust.
\draftnote{References for the images of all regions.}

\begin{table}[ht!]
\caption{The fields of nine regions, from the diffuse material below $10^{21.8}$\,cm$^{-2}$ (Appendix~\ref{app:fieldest}). $G_0$:
the fields that reproduce the 16th, 50th and 84th percentiles of the fitted temperature, as medians over the diffuse bins weighted
by their pixels. The field below the weakest one is marked with $^{<}0.98$. Too warm: the percentage of all pixels of the map
warmer than the strongest field of the grid allows.}
\label{tab:g0regions}
\centering
\begin{tabular}{lcc}
\hline\hline
\noalign{\smallskip}
Cloud & $G_0$ &\,\, \% \\
\noalign{\smallskip}
\hline
\noalign{\smallskip}
Taurus        & $^<$0.98          &\phantom{0}0.0\\
California    & 3.6\,\,5.4\,\,6.9 &\phantom{0}0.1\\
Perseus       & 3.2\,\,4.7\,\,6.1 &\phantom{0}0.1\\
Scorpius      & 6.2\,\,7.9\,\,10  &\phantom{0}0.0\\
Orion~A       & 5.9\,\,8.9\,\,28  &\phantom{0}2.0\\
Ophiuchus     & 14\,\,31\,\,55    &\phantom{0}0.7\\
Aquila        & 25\,\,33\,\,41    &\phantom{0}0.4\\
W3 W4 W5      & 54\,\,62\,\,73        &\phantom{0}1.3\\
Cygnus~X      & 200\,\,250\,\,320     & 21.8\\
\noalign{\smallskip}
\hline
\end{tabular}
\end{table}

A single field per cloud is an approximation. In Orion~A and Ophiuchus, the fields at the 16th and 84th percentiles differ by a
factor of four or five, because parts of these clouds are lit by nearby OB stars. For such clouds, the map can be divided into
regions with their own fields. In Cygnus~X, a fifth of the pixels are warmer than the strongest field of the grid allows, and the
correction is not calibrated for them.

\section{Why the sphere and the layer are the limits}
\label{app:shape}

\begin{figure*}[ht!]
\centering
\includegraphics[width=0.96\textwidth]{fig_angular.pdf}
\caption{\emph{Left}: The fraction of the heating radiation that survives the journey inward, against the direction it comes from, 
for a parcel at mid-depth of a plane-parallel layer of three total surface densities. Plotted against the cosine of the angle, the
angular mean of each curve is the area beneath it, drawn dashed; the single column that would produce that same mean intensity is
the effective attenuating column of Eq.~(\ref{eq:seff}). \emph{Right}: that column against the total surface density of the line of
sight, for the two limiting shapes and for the average of them that is adopted, with the shortest path for comparison.}
\label{fig:angular}
\end{figure*}

The cube in Fig.~\ref{fig:relation} shows that the sphere and the layer are indeed the limits. The cube has the symmetry of neither
of them. Its factor lies between them at every surface density, within $6$\% of the sphere and within $3$\% of the layer. The layer
is an extreme case because it is unbounded across the line of sight. For a parcel at the mid-plane, the path to the boundary is the
half-thickness divided by the cosine of the angle to the normal. This path diverges as the angle approaches $90\degr$, so a whole
band of directions near the plane of the layer delivers no radiation, however thin the layer is. At $10^{22}$\,cm$^{-2}$, the
directions within $30\degr$ of the normal deliver $64$\% of the mean intensity, and those within $45\degr$ deliver $93$\%. In a
bounded body, the path is finite in every direction, and no direction is lost. The sphere is the opposite extreme: a compact body
with a given central surface density blocks less of the sky than any other body with that surface density. Every real cloud lies
between the two limits, and the range between them is the full effect of the shape. This is why the method needs no assumption
about the shape of a cloud.

The dotted curve of Fig.~\ref{fig:relation} shows what happens if the angular average is replaced by the column along the shortest
path, which is one quarter of the surface density when averaged over the depth of a layer. Below $3{\times}10^{21}$\,cm$^{-2}$,
this simpler treatment agrees with the full calculation to $1$\%. Above $10^{22}$\,cm$^{-2}$ it diverges, and at
$10^{23}$\,cm$^{-2}$ it overestimates the correction by half. The shortest path gives too little shielding to the material near the
surfaces: its column falls to zero at the surface, while the oblique directions still shield the surface material. The surface
material is then too warm, and the required correction is overestimated.

\section{Corrections that read the fitted temperature}
\label{app:tcorr}

The reconstruction gives a dust temperature for every pixel, together with the surface density. If the temperature of a pixel
differs from the one that the relation predicts at its surface density, the temperature carries information that the surface
density does not. Four ways of using it were tested, all calibrated on the radiative transfer models. None was adopted, and the
reason is the same for all four.

The methods were: a table of the required factor against the surface density, the fitted temperature, and the contrast between the
images at $13.5{\arcsec}$ and $36.3{\arcsec}$; a smooth function of the surface density and of the offset of the fitted temperature
from the relation, fitted to the grid of radiation fields; a table computed from analytic lines of sight through structures
embedded in a diffuse medium; and a redistribution of the correction within each beam to follow the shape of the finest image. On
the models of the subdivision grid, they changed the median core mass by $2$ to $9$\% relative to the adopted correction. The third
method under-corrected the dense material in fields of $G_0 = 20$ to $50$ by a factor of $1.5$ to $1.7$. On an earlier simulated
sky, the first two over-corrected the flanks of filaments even more than the adopted correction, and the fourth copied the noise of
the finest image into the factor.

The reason is that a warm pixel is ambiguous. A pixel can be warmer than the relation predicts for two reasons. First, a stronger
radiation field may heat it. The radiative transfer models then say that it needs more correction. Second, it may lie on the flank
of a structure, near its surface, or its line of sight may also cross a warm cloud. It then needs less correction. The models used
for the calibration contain only the first case, and the simulated sky mostly the second. A correction that uses the temperature is
therefore right in one case and wrong in the other. The fitted temperature of a single pixel cannot distinguish the two cases. The
radiation field varies over parsecs, and the warmth of a flank varies over the width of a filament. The field is therefore
estimated from the temperature of the diffuse material of the whole map, as one number for a cloud (Appendix~\ref{app:fieldest}).
In a small field, such as the Aquila field of Sect.~\ref{sec:aquila}, only $12$\% of the pixels are diffuse enough, and most of
them lie in one corner, so the field must be taken from the larger map. None of these corrections changed the core mass by more
than the error due to the background (Appendix~\ref{sec:bgerr}), so the simplest one was adopted.

\end{document}
